\documentclass[10pt,a4paper,notitlepage,twocolumn,oneside]{article}

\usepackage[utf8]{inputenc}
\usepackage[a4paper,left=0.15in,right=0.15in,top=0.6in,bottom=0.6in]{geometry}
\usepackage{amsmath,amssymb,amsfonts}
\usepackage{hyperref}
\usepackage{graphicx}
\usepackage{endnotes}
\usepackage{fancyhdr}
\usepackage{titlesec}
\usepackage{setspace}
\usepackage{xcolor}
\usepackage{wrapfig}
\usepackage[font=footnotesize]{subcaption}
\usepackage{longtable,ltcaption}
\usepackage{xtab}
\usepackage{multirow}
\usepackage{array}
\usepackage{enumitem}
%\usepackage[american]{babel}
\usepackage{csquotes}
\usepackage{pdflscape}
\usepackage{comment}
\usepackage[backend=biber,
            style=apa,
            %style=numeric-comp,
            citestyle=authoryear-comp,
            %citestyle=numeric-comp,
            sorting=nyt,
            %sorting=none,
            %maxcitenames=2,
            %mincitenames=1
            ]{biblatex}

%% MACROS
\renewcommand{\footnoterule}{%
  \kern -3pt
  \hrule width \columnwidth height 0.1pt
  \kern 4.5pt
}
\newcommand{\orcid}[1]{\href{https://orcid.org/#1}{\includegraphics[width=1em]{orcid.pdf}}}
\newcommand{\scholar}[1]{\href{https://scholar.google.pt/citations?user=#1}{\includegraphics[width=1em]{scholar.pdf}}}

%% GENERAL DEFINITIONS
\DeclareUnicodeCharacter{0301}{\'{e}}
\urlstyle{same}
\DeclareCiteCommand{\citetitle}
  {\boolfalse{citetracker}%
   \boolfalse{pagetracker}%
   \usebibmacro{prenote}}
  {\ifciteindex
     {\indexfield{indextitle}}
     {}%
   \printtext[bibhyperref]{\printfield[citetitle]{labeltitle}}}
  {\multicitedelim}
  {\usebibmacro{postnote}}
\newcolumntype{P}[1]{>{\centering\arraybackslash}p{#1}}
\renewcommand{\arraystretch}{1.25}
\setlength{\tabcolsep}{3pt}

%% HEADINGS
\pagestyle{myheadings}
%\fancyhf{}
\markboth{Ricardo B. Sousa \textit{et al.}}{A Systematic Literature Review on Long-Term Localization and Mapping for Mobile Robots}

%% BIBLATEX
\addbibresource{references.bib}



%%%%%%%%%%%%%%%%%%%% ARTICLE INFORMATION

\title{A Systematic Literature Review on Long-Term Localization and Mapping\\for Mobile Robots%
\thanks{%
This work is financed by National Funds through the Portuguese funding agency, FCT -- Fundação para a Ciência e a Tecnologia, within project LA/P/0063/2020. This work is also financed by National Funds through the Portuguese funding agency, FCT -- Fundação para a Ciência e a Tecnologia, within scholarship 2021.04591.BD.

\vspace{1em}
\noindent Ricardo B. Sousa\textsuperscript{1, 2}

\noindent \href{mailto:up201503004@edu.fe.up.pt}{up201503004@edu.fe.up.pt}

\vspace{0.5em}
\noindent H\'{e}ber M. Sobreira\textsuperscript{2}

\noindent \href{mailto:heber.m.sobreira@inesctec.pt}{heber.m.sobreira@inesctec.pt}

\vspace{0.5em}
\noindent Ant\'{o}nio Paulo Moreira\textsuperscript{1, 2}

\noindent \href{mailto:amoreira@fe.up.pt}{amoreira@fe.up.pt}

\vspace{1em}
\noindent \textsuperscript{1} Faculty of Engineering of the University of Porto, Electrical Engineering Department, Porto, 4200-465 Porto, Portugal

\noindent \textsuperscript{2} INESC TEC -- Institute for Systems and Computer Engineering, Technology and Science, CRIIS -- Centre for Robotics in Industry and Intelligent Systems, Porto, 4200-465 Porto, Portugal
}}
\author{%
Ricardo B. Sousa
\href{https://orcid.org/0000-0003-4537-5095}{\includegraphics[width=1em]{orcid.pdf}}%
\and%
H\'{e}ber M. Sobreira
\href{https://orcid.org/0000-0002-8055-1093}{\includegraphics[width=1em]{orcid.pdf}}%
\and%
Ant\'{o}nio Paulo Moreira
\href{https://orcid.org/0000-0001-8573-3147}{\includegraphics[width=1em]{orcid.pdf}}%
}
%\date{\today}
\date{October 31, 2022}



%%%%%%%%%%%%%%%%%%%% ARTICLE BODY

\begin{document}

\maketitle

\section*{Abstract}

Long-term operation of autonomous robots creates new challenges to the Simultaneous Localization and Mapping (SLAM). Varying conditions of the vehicle's surroundings, such as appearance variations (lighting, daytime, weather, or seasonal) or reconfigurations of the environment, are a challenge for SLAM algorithms to adapt to new changes while preserving old states. When also operating for long periods and trajectory lengths, the map should readjust to environment changes but not grow indefinitely, where the map size should be dependent only on the explored environment area.
Long-term SLAM intends to overcome the challenges associated with lifelong autonomy and improve the robustness of autonomous systems. Although several studies review SLAM algorithms, none of them focus on lifelong autonomy.
Thus, this paper presents a systematic literature review on long-term localization and mapping following the Preferred Reporting Items for Systematic reviews and Meta-Analysis (PRISMA) guidelines. The review analyzes 142 works covering appearance invariance, modeling the environment dynamics, map size management, multi-session, and computational issues including parallel computing and timming efficiency. The analysis also focus on the experimental data and evaluation metrics commonly used to assess long-term autonomy.
Moreover, an overview over the bibliographic data of the 142 records provides analysis in terms of keywords and authorship co-occurrence to identify the terms more used in long-term SLAM and research networks between authors, respectively.
Future studies can update this paper thanks to the systematic methodology presented in the review and the public GitHub repository with all the documentation and scripts used during the review process.

\hfill\break

\noindent\textbf{Keywords:} simultaneous localization and mapping (SLAM), lifelong SLAM, long-term autonomy, mobile robots.

\section{Introduction}
\label{sec:intro}

An autonomous mobile robot requires a representation of its surroundings to localize itself relative to the environment.
Simultaneous Localization and Mapping (SLAM) addresses this problem by incorporating the robot state estimation (pose and possibly other state variables) concurrently with the mapping process. This process builds a representation of the environment perceived by the robot originating a map incrementally built when exploring unknown areas or refined on passages through known locations.

In a static scene, the robot would only need to map once because it would be always consistent with the environment.
However, autonomous systems deployed in industrial locations, outdoor environments, or even service-oriented applications such as shopping centers or homes deal with moving elements in the scene (humans, objects), environment reconfiguration (machines moving places), and appearance variations (lighting, weather, seasonal, or daytime changes).
These varying conditions are a challenge for the SLAM system, where the system should decide how (consider the most current state or only the most static changes?) and when to update the map (when the variations occur or after a certain time?).
This challenge is also known as the stability-plasticity dilemma, where the long-term localization and mapping should both adapt to new environment changes and preserve old states over time~\parencite{biber-duckett:2009:0278364908096286}.

Furthermore, the map would grow indefinitely when gathering new information from the environment. This ever-growing problem poses another challenge for the SLAM system in the long-term due to the limited computational resources of the mapping vehicle. Indeed, the map size should be dependent on the environment area and not on the trajectory length of the mapping process nor operation time, only growing when the robot would explore unknown locations~\parencite{kretzschmar-stachniss:2012:0278364912455072}.

Although several studies overview SLAM literature, only a subset of those studies mentions long-term challenges of performing SLAM.
\cite{purpose:study:cadena:2016} has a brief survey on the robustness and scalability of autonomous systems focused on loop closure validation, dynamic environments, pose graph sparsification, and parallel and distributed computing for metric and semantic SLAM.
In contrast, \cite{purpose:study:lowry:2016} limits the study to vision-based topological SLAM discussing also the challenge of varying conditions.
While \cite{purpose:study:bresson:2017} overviews autonomous driving trends in terms of scalability, map updatability, and dynamicity, the survey limits the discussion to algorithms that have both odometry and mapping modules, excluding localization-only works. Also, \cite{purpose:study:bresson:2017} focuses more on the modules of the SLAM (relocalization, localization against a map).
\cite{purpose:study:kunze:2018} gives a brief overview of AI-related works for robustness to appearance changes and learning dynamics of the environment, discussing areas in which artificial intelligence (AI) enables long-term operation of autonomous systems.
As for \cite{purpose:study:zaffar:2018}, the study evaluates the long-term autonomy of sensors such as monocular and stereo cameras, and LiDAR.
However, to the best of the authors knowledge, none of the existing studies overviews the trends for dealing with long-term challenges in SLAM. Also, the studies that overview some of the challenges of lifelong SLAM do not clarify the process for identifying the cited works, not allowing other researchers to repeat the identification process of records when searching in bibliographic databases.

Therefore, this paper presents a systematic literature review on long-term localization and mapping following the Preferred Reporting Items for Systematic reviews and Meta-Analysis~(PRISMA)~\parencite{methodology:prisma} statement.
The systematic method followed in this review allows the replication of the results by other researchers and leads to the inclusion of 142 works for discussion and analysis.
Also, this paper makes available a public GitHub repository\footnote{\url{https://github.com/sousarbarb/slr-ltlm-mr}} with all the documentation and scripts used during the process of systematic revision of the literature.
The main contributions of this paper relative to the existing studies on SLAM are the following ones:

\begin{itemize}[nosep]
\item discussion on methodologies and trends focused on appearance invariance, dynamic elements, map sparsification and multi-session techniques, and other computational concerns;
\item comparative analysis on the public datasets and experimental data used in the included works in terms of environment conditions, sensorization, and distance and time properties;
\item presentation of common evaluation metrics used by the included works in the experimental results.
\end{itemize}

This review does not intend to review the fundamentals of SLAM nor the main formulations. The reader should refer to \cite{purpose:study:durrant-whyte-bailey:2006:1} and \cite{purpose:study:durrant-whyte-bailey:2006:2} for the Extended Kalman Filter (EKF) and particle filter probabilistic formulations, and to \cite{purpose:study:grisetti:2010} for the pose graph formulation of SLAM.





\subsection{Paper organization}

The rest of this review is organized as follows.
Section~\ref{sec:purpose} discusses the limitations of existent studies and reviews and presents the purpose and motivations of this paper.
Section~\ref{sec:methodology} explains the methodology followed in this review to search and select the included records, and the data extraction process for synthesis and analysis. In Appendix~\ref{a2:data-extraction}, Table~\ref{tab:a2:data-extraction} presents the data extraction results of the included records.
Next, Section~\ref{sec:overview} analyzes the bibliographic information of the 142 included works in this review in terms of the identification results of each data source considered in the methodology, keywords co-occurrence and co-authorship relations, the year of publication, and the publication venue.
Section~\ref{sec:discussion} discusses the methodologies found in the included records related to long-term localization and mapping and analyzes the experimental data and evaluation metrics used in the experiments by the authors.
Then, Section~\ref{sec:future} outlines challenges and future directions.
Section~\ref{sec:limitations} discusses possible limitations of this study.
Lastly, Section~\ref{sec:conclusions} presents the conclusions.

\section{Purpose of the study}
\label{sec:purpose}

\textcolor{red}{PUT INTRO HERE!!!}





\subsection{Limitations of current studies}



\begin{table}[h]
  \caption[Existent Literature Reviews, Surveys, and Tutorials on SLAM.]{Existent Literature Reviews, Surveys, and Tutorials on SLAM.}
  \label{tab:purpose:current-literature}
  \centering
  {\scriptsize
  \begin{tabular}{c p{0.43\columnwidth} p{0.43\columnwidth}}

\hline
\textbf{Year} & \textbf{Topic} & \textbf{Reference}\\
\hline
2006 & Probabilistic formulations (EKF, particle filter)
     & \cite{purpose:study:durrant-whyte-bailey:2006:1}, \cite{purpose:study:durrant-whyte-bailey:2006:2}\\
\hline
2008 & Probabilistic and pose graph formulations
     & \cite{purpose:study:thrun:2008}\\
\hline
2010 & GraphSLAM
     & \cite{purpose:study:grisetti:2010}\\
\hline
2011 & Observability, convergence, consistency (feature-based SLAM)
     & \cite{purpose:study:dissanayake:2011}\\
\hline
2012 & Visual odometry
     & \cite{purpose:study:scaramuzza-fraundorfer:2011:1}, \cite{purpose:study:scaramuzza-fraundorfer:2012:2}\\
\hline
2014 & Underwater navigation and localization
     & \cite{purpose:study:paull:2014}\\
\hline
2015 & Visual SLAM
     & \cite{purpose:study:yousif:2015}\\
\hline
2016 & Observability, convergence, consistency (feature and graph-based SLAM)
     & \cite{purpose:study:huang-dissanayake:2016}\\
     & Multi-robot SLAM
     & \cite{purpose:study:saeedi:2016}\\
     & Visual place recognition
     & \cite{purpose:study:lowry:2016}\\
     & SLAM literature overview
     & \cite{purpose:study:cadena:2016}\\
\hline
2017 & Autonomous vehicles
     & \cite{purpose:study:bresson:2017}\\
\cline{1-1}
\cline{3-3}
2018 &
     & \cite{purpose:study:kuutti:2018}\\
\hline
2018 & AI for long-term autonomy
     & \cite{purpose:study:kunze:2018}\\
     & Long-term sensorization
     & \cite{purpose:study:zaffar:2018}\\
\hline
2020 & Deep learning
     & \cite{purpose:study:fayyad:2020}\\
     & Multi-robot search and rescue
     & \cite{purpose:study:queralta:2020}\\
\hline
2021 & Self-driving vehicles
     & \cite{purpose:study:badue:2021}\\
\hline
2022 & Underground navigation
     & \cite{purpose:study:ebadi}\\
\hline

  \end{tabular}}
\end{table}





\subsection{Motivations and goals}

Research question: What is the current state of the art of long-term localization and mapping using mobile robots?

Goals of this review:

\begin{itemize}[nosep]
\item which are the main strategies for accomplishing long-term operations with mobile robots;
\item how to deal with varying conditions of the environment;
\item how do autonomous robots deal with the dynamics of the environment;
\item which are the main strategies to deal with the limited computational resources of a mobile robot on long-term operations;
\item how the methods evaluate their results;
\item which are the public datasets more used for evaluating long-term localization and mapping.
\end{itemize}

PICO framework (Population--Intervention--Comparison--Outcome) helps to frame the research questions of this systematic review into a more structured framework:

\begin{itemize}[nosep]
\item \textbf{Population:} mobile robots;
\item \textbf{Intervention:} localization, mapping, SLAM;
\item \textbf{Comparison:} \textit{not applicable to this study};
\item \textbf{Outcome:} long-term operation, lifelong autonomy, robust.
%\item \textbf{Context:} continuous operation, service robots, industrial environments.
\end{itemize}


\section{Methodology}
\label{sec:methodology}

A systematic literature review uses explicit, rigorous, and reproducible systematic methods to synthesize the findings of studies related to a particular research question, topic area, or phenomenon of interest. This type of review assures the quality and trustworthiness of the review's findings by presenting a complete, organized, and summarized analysis of all works considered while allowing others to replicate or update the reviews. The most common standard for performing a systematic review is the PRISMA~\parencite{methodology:prisma} statement. Although the PRISMA statement has been designed originally for evaluating the effects of health interventions, the checklist items of the methodology are general and applicable to other subject areas. Thus, the methodology used in this systematic review follows the PRISMA~\parencite{methodology:prisma} guidelines.

This section presents the detailed methodology used in this study. First, the eligibility criteria decide which studies to include in the review. Next, the search strategy details the information sources considered in the review and the base string and search fields used for inquiring these sources. Furthermore, the selection process focuses on describing its stages and the quality evaluation criteria used to select works for the synthesis and analysis phase of the review. Lastly, the data extraction process details the relevant data collected for synthesis and analysis.
Parsifal \parencite{methodology:parsifal} is the online tool used to support the literature review in designing the methodology protocol, removing duplicates, screening and selecting works including their quality assessment. Additional documentation and scripts developed within the scope of this review related to removing duplicates, checking and processing the bibliographic references, and data extraction are available in the public GitHub repository.

\subsection{Eligibility criteria}
\label{sec:methodology:eligible}

Table~\ref{tab:methodology:exclusion-criteria} presents the exclusion criteria used to determine the eligible studies for the selection process. These eligibility criteria focus mainly on the type of paper and availability. The index criterion rejects all publications not indexed in a scientific publication venue. This rejection guarantees that the eligible works were peer-reviewed by the scientific community. Also, the exclusion criteria reject short papers and gray, secondary, and tertiary literature. Short papers do not usually present a detailed methodology of their scientific contribution. As for only considering primary literature in the review, this criterion increases the relevance of search results by favoring original articles and simultaneously guaranteeing peer-revision of the works. In terms of language, only considering studies with English full-texts increases the scope and visibility of the review. Similarly, the eligibility criteria reject studies not available in digital libraries for reproducibility and accessibility reasons.

\begin{table}[h]
  \renewcommand{\arraystretch}{1.25}
  \setlength{\tabcolsep}{3pt}
  \caption[Exclusion criteria for the selection process.]{Exclusion criteria for the selection process.}
  \label{tab:methodology:exclusion-criteria}
  \centering
  {\scriptsize
  \begin{tabular}{c m{0.28\columnwidth} m{0.6\columnwidth}}

\hline
\textbf{E\#} & \textbf{Criteria} & \textbf{Statement}\\
\hline
E1 &
Index &
Papers not indexed in a scientific publication venue\\
\hline
E2 &
Language &
Full-text of the papers not published in English\\
\hline
E3 &
Subject Area &
Papers not classified in the databases as Computer Science, Engineering, Mathematics, or Multidisciplinary\\
\hline
E4 &
Short Papers &
Papers classified as short papers accordingly to the publication venue\\
\hline
E5 &
Gray, Secondary, and Tertiary Literature &
Books, preprints, reports, reviews, thesis, ...\\
\hline
E6 &
Availability &
Full-text of the papers not available in digital libraries\\
\hline
E7 &
Dataset &
Papers that focus only on data collection\\
\hline
E8 &
Coverage &
Papers using only odometry for localization\\
\hline
E9 &
Scope &
Papers that focus on different and not related subjects\\
\hline

  \end{tabular}}
\end{table}

Another exclusion criterion considered in the review is relative to the studies' categorization of their subject areas by bibliographic databases. The ones considered in the review are Computer Science, Engineering, Mathematics, or Multidisciplinary areas. In the list provided by the Clarivate's Journal Citation Reports\footnote{\url{https://jcr.clarivate.com/jcr/browse-categories}}, these four subject areas include the artificial intelligence, interdisciplinary applications, electrical and computers engineering, robotics, and applied mathematics categories, among others. These categories are intrinsically related to the localization and mapping problem for long-term operation of mobile robots.

The final three criteria presented in Table~\ref{tab:methodology:exclusion-criteria} focus on the scientific contribution of the studies. The dataset criterion rejects all works that focus only on sharing a data collection. Although these works are important for the evolution of localization and mapping algorithms in providing a benchmark for comparison and reference purposes, their scientific contribution is not directly comparable to research articles. Odometry-only approaches are unusable over long distances invalidating their use for long-term operations with mobile robots. As for the scope criterion, this review does not consider as eligible papers not related to long-term localization and mapping that do not fall under the other exclusion criteria.

\subsection{Search strategy}
\label{sec:methodology:search}

The search phase consists of identifying the data sources that could be relevant for this literature review, and defining the base string and which search fields considered to obtain the results for the review. \citetitle{methodology:search:db:wos} and \citetitle{methodology:search:db:scopus} are traditionally the two most widely used bibliographic databases. However, previous studies demonstrate that different databases differ significantly in their scientific coverage~\parencite{methodology:search:db:coverage:scopus-wos,methodology:search:db:coverage:dim-scopus-wos}. Thus, the data sources considered in this review are the following ones: \citetitle{methodology:search:db:acm}, \citetitle{methodology:search:db:dimensions}, \citetitle{methodology:search:db:ieee-xplore}, \citetitle{methodology:search:db:inspec}, \citetitle{methodology:search:db:scopus}, and \citetitle{methodology:search:db:wos}.

Moreover, May 17, 2022, is the date of the last full inquiry. Future reviews on the topic of this study should consider this final date as theirs initial one. As for inquiring the data sources, the base string used is the following one:

\vspace{1em}

\noindent\begin{center}
\begin{tabular}{l l}
\texttt{(robot* OR vehicle*)}                     & \texttt{AND}\\
\texttt{((locali* AND map*) OR "slam")}           & \texttt{AND}\\
\texttt{("long term" OR "life long" OR lifelong)}\\
\end{tabular}
\end{center}

\vspace{1em}

The first terms, \texttt{robot* OR vehicle*}, attempt to focus the search results to the desired population. These two terms have multiple synonyms within the scope of autonomous mobile robots: mobile robots, autonomous vehicles, robotics, agricultural robots, intelligent robots, service robots, unmanned aerial/ground/underwater vehicles, among other terms. Therefore, by adding the asterisk to the end of the terms robot and vehicle (\texttt{robot*} and \texttt{vehicle*}, respectively), and by only considering the terms with asterisk in the inquiry, all the synonyms are covered for the desired population. Given the incompatibility of the \citetitle{methodology:search:db:dimensions} database with wildcards (e.g., using the asterisk), the first part of the base string becomes as follows when searching in this database: \texttt{robot OR robots OR robotics OR vehicle OR vehicles}.

The next part of the query focus on the intervention side of the systematic review. Given the interest of this review on searching for localization and mapping algorithms, \texttt{locali*} and \texttt{map*} summarize all the synonyms for the localization and mapping terms, respectively. For example, \texttt{locali*} not only is agnostic to the US versus UK spelling differences (localization vs localisation, respectively) but also resumes several synonyms: localization, localize, or localizing. The term \texttt{map*} also attempts to cover its respective synonyms such as map, maps, or mapping.
Also, the acronym \texttt{"slam"} is another alternative to search for localization and mapping algorithms. Even though its definition is compatible with \texttt{locali* AND map*}, some authors only refer to SLAM.
Similarly to the inquiry's first part, the second one becomes as follows for searching in \citetitle{methodology:search:db:dimensions}: \texttt{((localize OR localization OR localizing OR localise OR localisation OR localising) AND (map OR maps OR mapping)) OR "slam"}.

As for \texttt{"long term" OR "life long" OR lifelong}, this part of the base string is relative to the outcome of the PICO framework, presented in Section~\ref{sec:purpose}. The reason for having both \texttt{"life long"} and \texttt{lifelong} terms is the existing confusion in which term is grammatically the correct one.

Furthermore, the Title, Abstract, and Keywords are the fields considered for obtaining the search results. The third one includes the author keywords, the indexed terms by the databases, and the uncontrolled ones if they are available. The selection of these search fields for this review improves the relevance of the results compared to using all fields and the full text by focusing the search on the summary items of the works. Indeed, the main contributions of scientific works should be summarized in at least the title, abstract, or the author keywords. The indexed terms also help in obtaining records only related to the base string used in this review.
However, not all data sources have available the search fields considered in the review or some of them require an adaptation when performing the search. 
Although the \citetitle{methodology:search:db:acm} allows searching within multiple search fields, including the ones considered in this review, the advanced search query on this library sets by default an AND operator between the different fields. This setting must be changed manually in the query syntax to the desired OR operator. Also, there are two options to search items in the \citetitle{methodology:search:db:acm}: \textit{The ACM Full-Text Collection} and \textit{The ACM Guide to Computing Literature}. Given that the latter includes all the content from the former, the identification process in this source performs the search using \textit{The ACM Guide to Computing Literature} option.
Other than searching in the publications' full data, \citetitle{methodology:search:db:dimensions} only has the title and abstract search fields compatible with this review. Given the limitation of \citetitle{methodology:search:db:ieee-xplore} to 7 wildcards, the search results of this digital library using the base string for the inquiry are the grouping of different searches considering only a search field at a time, importing each search results to Parsifal and removing the duplicates. As for \citetitle{methodology:search:db:inspec}, \citetitle{methodology:search:db:scopus}, and \citetitle{methodology:search:db:wos}, these databases have available all the search fields considered in the review.

In terms of the publication date, this review does not restrict it to avoid ignoring important works and to improve the discussion. Indeed, to best of the authors knowledge, there is not available a systematic review on long-term localization and mapping for mobile robots to provide an initial date for rejecting older publications. Even though the number of publications per year could indicate an initial date on when the topic gained relevance, the date filtering could still reject important works.

\subsection{Selection process}
\label{sec:methodology:selection}

The selection process of this review summarized in Figure~\ref{fig:methodology:prisma-flow} has three phases: identification, screening, and quality assessment. The first phase consists of inquiring each data source discussed previously with the base string and adapting it if needed. The second phase requires screening the papers. In this review, screening is equivalent to reading the publications' title and abstract and deciding whether the study is eligible or not based on the exclusion criteria. Then, a set of evaluation criteria assesses the quality of the eligible records. The records obtained after the three phases of the selection process are for the data extraction phase.

\begin{figure}[h]
  \centering
  \includegraphics[width=\columnwidth]{figures/selection.pdf}
  \caption{PRISMA flow diagram for the selection process.}
  \label{fig:methodology:prisma-flow}
\end{figure}

\subsubsection{Identification}

In the identification phase of this review, the search strategy is applied to all data sources. \citetitle{methodology:search:db:acm}, \citetitle{methodology:search:db:dimensions}, \citetitle{methodology:search:db:inspec}, \citetitle{methodology:search:db:scopus}, and \citetitle{methodology:search:db:wos} data sources only require a single inquiry to obtain the search results.
Given the limitation of the \citetitle{methodology:search:db:ieee-xplore} for using wildcards mentioned in Section~\ref{sec:methodology:search}, the number of records for this source presented in Figure~\ref{fig:methodology:prisma-flow} represents the results of 7 inquiries (using the fields title, abstract, author keywords, IEEE terms, INSPEC controlled terms, and the INSPEC uncontrolled ones, respectively) after removing the duplicates with the support of Parsifal. Although the total number of search results found is 2160, Parsifal is used to remove duplicates from different data sources, excluding 1339 records. Following the duplicates removal, the exclusion criteria defined in Section~\ref{sec:methodology:search} exclude 232 works from the review. This exclusion is possible due to \citetitle{methodology:search:db:inspec}, \citetitle{methodology:search:db:scopus}, or \citetitle{methodology:search:db:wos} having filters related to the publication's type, subject area, and language.

The works excluded from the search results also include the ones that do not meet the exclusion criteria E4 and E7. For the first one, a Python script available in the GitHub repository of this review searches studies with a number of pages lower or equal to 4. Even though short papers have a maximum number of 3 pages, the papers with 4 pages do not usually present a detailed methodology.
As for the E7 exclusion criterion, some works are possible to remove from the review by searching in their title for the term ``dataset''.
All excluded articles of this review are double-checked to certify if the exclusion criteria are correctly applied. For example, articles published in the Remote Sensing journal from MDPI do not meet the E3 criterion. Indeed, the Journal Citations Reports from Clarivate classifies it by the following categories: Remote Sensing, Geosciences Multidisciplinary, Environmental Sciences, and Imaging Science \& Photographic Technology. However, most search results from this journal found in the identification phase are directly related to the topic of this review and the respective subject areas. Thus, in these cases and in other ones related to the remaining exclusion criteria, the decision is reverted to consider the initially rejected studies for the next phase of the review.

\subsubsection{Screening}

Next, the screening phase in this review consists of reading the title and abstract of the publications and rejecting the ones that meet the exclusion criteria. However, the initially rejected papers have another assessment for validating the exclusion. The analysis of the results and conclusions of these publications considering the exclusion criteria either confirms the exclusion decision or reverses it to eligible works for quality assessment. As a result of the screening phase, 178 studies are rejected from the initial identified 589 works. The duplicate records found in screening and removed manually are due to titles with invalid characters originated by exporting the search results from the \citetitle{methodology:search:db:dimensions}~database.

\subsubsection{Quality assessment}

The quality evaluation in this review of the selected works from screening follows the 8 Quality Evaluation (QE) criteria presented in Table~\ref{tab:methodology:quality-assessment}. All of them are subjective criteria derived from the analysis of the eligible works. The score column establishes the possible values for the QE criteria, in which the minimum, intermediate, and maximum values correspond to none, partial, and full compliance, respectively.
Furthermore, QE1, QE2, QE4, and QE8 focus on the details provided in the papers, specifically, if the discussion of the related work, the proposed methodology, the experimental setup, and the results are detailed and thoroughly analyzed in the publication, respectively.
The possible scores for QE3 are twice the value of QE1, QE2, QE4, and QE8 due to this criterion being directly related to the topic of the review. A work focusing on both localization and mapping problems will have a score of 2.0 (full compliance). If the study only focuses on one of these problems or none of them, the scores will be 1.0 or 0.0, i.e., partial or no compliance, respectively.
QE5 evaluates the long-term results of the eligible studies and is either 2.0 (full) or 0.0 (no compliance). This criterion has the same range as QE3 for similar reasons, given the focus of this review on long-term localization and mapping algorithms.
The definition of long-term experiments for assigning full compliance in QE5 is the following one: dynamic changing environments (e.g., dynamic elements or semi-static ones), increasing environments or feature maps in terms of their size, redundant data removal, or varying conditions (e.g., different seasons of the year or lighting conditions).
QE6 and QE7 can only be 1.0 or 0.0. The former criterion intends to highlight works that compare themselves to the state of the art and/or ground-truth data. The latter emphasizes the importance of having available either the implementation of the proposed methodology or the data used in the experiments for other works to be able to compare the proposed methodologies.
Lastly, considering the possible scores for the QE criteria in Table~\ref{tab:methodology:quality-assessment}, each work can only have a  maximum score of 10.0.

\begin{table}[h]
  \renewcommand{\arraystretch}{1.25}
  \setlength{\tabcolsep}{3pt}
  \caption[Quality evaluation criteria and score range.]{Quality evaluation criteria and score range.}
  \label{tab:methodology:quality-assessment}
  \centering
  {\scriptsize
  \begin{tabular}{c m{0.68\columnwidth} c}

\hline
\textbf{QE\#} & \textbf{Criteria} & \textbf{Score}\\
\hline
QE1 &
Does the paper have an updated state of the art on long-term localization and mapping? &
\{0.0, 0.5, 1.0\}\\
\hline
QE2 &
Is the methodology appropriate and detailed? &
\{0.0, 0.5, 1.0\}\\
\hline
QE3 &
Does the methodology consider both localization and mapping problems? &
\{0.0, 1.0, 2.0\}\\
\hline
QE4 &
Is the hardware and/or software used in the experiments detailed? &
\{0.0, 0.5, 1.0\}\\
\hline
QE5 &
Does the paper presents any kind of long-term experimental results? &
\{0.0, 2.0\}\\
\hline
QE6 &
Does the paper presents comparative results with other methods and/or ground-truth data? &
\{0.0, 1.0\}\\
\hline
QE7 &
Does the work's implementation and/or the data used in the experiments are publicly available? &
\{0.0, 1.0\}\\
\hline
QE8 &
Is the discussion of the results and conclusions appropriate and detailed? &
\{0.0, 0.5, 1.0\}\\
\hline

  \end{tabular}}
\end{table}

After evaluating the 411 eligible works accordingly to the previously discussed QE criteria (the scores of each record are available in the GitHub repository), the first conclusion of the authors is that works with a non-detailed or not appropriate methodology, results' discussion, or conclusions should not be included in the review. Another conclusion is relative to rejecting works that do not consider either localization or mapping problems, or do not present any long-term experimental results, given the focus of this review on the long-term localization and mapping problem for mobile robots. Furthermore, the quality assessment phase should consider a cut-off score to filter works with low quality scores. Consequently, the assessment phase considers the following two reasons to reject a record:

\begin{enumerate}[nosep]
\item QE2, QE3, QE5, QE8: reject works with a 0.0 (no compliance) score;
\item cut-off score: reject works with a score lower or equal to 7.5/10.0.
\end{enumerate}

The distribution of the evaluation scores and the QE criteria itself justify the selection of a 7.5/10.0 cut-off score. 
Figure~\ref{fig:methodology:qe} illustrates the scores distribution for all eligible works versus the scores of the ones that pass the first criterion defined previously for the QE phase (related to the compliance on the QE2, QE3, QE5, and QE8 criteria). The assessment of this criterion rejects 116 records (28\%) of the 411 eligible works (see Figure~\ref{fig:methodology:prisma-flow}). 
Even though the distribution of the evaluation scores changes significantly in the range of scores lower or equal to 7.5/10.0, as observed between Figures~\ref{fig:methodology:qe:qe_wo-r1} and~\ref{fig:methodology:qe:qe}, only one work with a score higher than 7.5 is rejected due to not having a detailed and appropriate discussion of the results. This result indicates that interesting works are associated with high scores, as intended when using a quality assessment methodology, while also suggests that the range between 8.0 and 10.0 have the most interesting and quality works compatible with the focus of this review on long-term localization and mapping. Although only assessing the eligible works would seem to lead to the same results in terms of records included in the review, the rejection criterion on QE2/3/5/8 prevents outliers related to the quality assessment.
From the remaining 295 eligible works, cut-off scores from 7.5 up to 8.5 have the following corresponding rejection rates:

\begin{itemize}[nosep]
\item 7.5/10.0 \phantom{$\xrightarrow{\text{reject}}$} 120 records (40.7\%) \phantom{$\xrightarrow{\text{include}}$} 175 records
\item 8.0/10.0 $\xrightarrow{\text{reject}}$ 160 records (54.2\%) $\xrightarrow{\text{include}}$ 135 records
\item 8.5/10.0 \phantom{$\xrightarrow{\text{reject}}$} 203 records (68.8\%) \phantom{$\xrightarrow{\text{include}}$} 92 records
\end{itemize}

\begin{figure}[h]
  \centering
  \subfloat[][]{\includegraphics[width=\columnwidth]{figures/qe.png}%
  \label{fig:methodology:qe:qe}}
  \linebreak
  \subfloat[][]{\includegraphics[width=\columnwidth]{figures/qe_wo-r1.png}%
  \label{fig:methodology:qe:qe_wo-r1}}
  \caption{Distribution of the quality evaluation scores obtained from assessing the eligible works considered in the review: (a) all eligible works; (b) works that pass the rejection criterion during the QE assessment related to QE2/3/5/8 = 0.0 (no compliance).}
  \label{fig:methodology:qe}
\end{figure}

The 8.5 cut-off score would not be suitable because methods that focus only on localization or mapping, or not having either the implementation or the experimental data publicly available would be obligated to have maximum scores in the other criteria to be included in the review. In these cases, a work would have a maximum score of 9.0 due to partial compliance on QE3 or no compliance on the QE7 criteria. Likewise, a cut-off score of 8.0 would only leave a margin for having a single partial compliance on QE1, QE2, QE4 or QE8 criteria in similar cases, even though it would reject 160/295 (54\%) records. Therefore, the 7.5/10.0 cut-off score is more appropriate for the quality assessment phase in this review by leaving margin for works to have partial compliance in more than one criterion. Indeed, this cut-off score allows an article with no public data and/or implementation (e.g., due to confidentiality agreements) to have up to four criteria with partial compliance, depending on the criterion's maximum score or if the work has available the experiments data and/or implementation. Another example is articles that only focus on localization or mapping. In these cases, the work could have no public implementation, even though requiring a maximum score on all other criteria, or, if the work has public data or implementation available, two other criteria could have partial compliance.

Overall, as illustrated in Figure~\ref{fig:methodology:prisma-flow}, the quality assessment of the 411 eligible works considering the two rejection criteria previously mentioned leads to rejecting a total of 236 (57\%) records. As a result, the remaining 175 records will be analyzed for data extraction.

\subsection{Data extraction}
\label{sec:methodology:data}

The data extraction process analyzes the records selected after the quality assessment phase and extracts information from these works. In the scope of this review, the Data Extraction (DE) items required for each record are the following ones:

\begin{itemize}[nosep]
\item $\left[\textbf{DE1}\right]$ \textbf{Long-term considerations} -- long-term factors the works consider in their proposed approach and experiments. Considering the knowledge obtained in the previous phases of this review's methodology, the authors considered the following factors for categorizing the included works:
  \begin{itemize}[nosep]
  \item appearance: varying conditions, appearance changes;
  \item dynamics: environment dynamics, dynamic elements;
  \item sparsity: map pruning, redundant data removal;
  \item multi-session: map management;
  \item computational: memory management, efficiency.
  \end{itemize}
\item $\left[\textbf{DE2}\right]$ \textbf{Localization} -- how the robot localizes itself and the type of localizer;
\item $\left[\textbf{DE3}\right]$ \textbf{Mapping} -- type of the map;
\item $\left[\textbf{DE4}\right]$ \textbf{Multi-robot} -- if the proposed methodologies consider multi-robot systems;
\item $\left[\textbf{DE5}\right]$ \textbf{Execution mode} -- offline, online, if requires both, or if no information on this item;
\item $\left[\textbf{DE6}\right]$ \textbf{Environment and domain} -- type of environment (indoor, outdoor) and domains (air, ground, water) tested with the proposed methodologies;
\item $\left[\textbf{DE7}\right]$ \textbf{Sensory setup} -- which sensors considered in the methodologies;
\item $\left[\textbf{DE8}\right]$ \textbf{Non-public experiments} -- if the authors performed experiments or tests with non-public data;
\item $\left[\textbf{DE9}\right]$ \textbf{Ground-truth} -- how ground-truth for non-public data is obtained or its type, if available;
\item $\left[\textbf{DE10}\right]$ \textbf{Distance and time characteristics} -- relative to the non-public experiments if available, as follows:
  \begin{itemize}[nosep]
  \item total distance (km) of the non-public experiments;
  \item path (km), in the case of repetitive paths;
  \item total time (h) in terms of continuous operation;
  \item time interval (day/week/month/year, or d/w/m/y) between the first and the last run.
  \end{itemize}
\item $\left[\textbf{DE11}\right]$ \textbf{Datasets} -- if and which public datasets are used in the experiments;
\item $\left[\textbf{DE12}\right]$ \textbf{Evaluation metrics} -- which metrics are used for evaluation.
\end{itemize}

In Section~\ref{sec:discussion:experiments}, a comparison table of the public datasets identified by the DE11 will contain the sensory setup, ground-truth data availability from the datasets, and the distance and time characteristics, similar to the data extraction items for non-public data, among other aspects. As a result, the distinction between public and non-public data availability represented in DE8, DE9, and DE10 allows to understand the distance and time characteristics of non-public data independently from the public datasets.

Although the data extraction phase in a systematic literature review usually does not remove any records, 33 of the analyzed 175 works have extended versions of the proposed methodologies, more detailed ones, or equivalent methods applied in different conditions.
Thus, these records are not included in the review to improve the discussion section in terms of singularity and originality of proposed approaches for the long-term localization and mapping problem.
The extracted information helped identifying the corresponding extended and more complete versions of these works.
A document containing the association of the removed versions to the records included in the review is available in the public GitHub repository, including their bibliographic references. 
Consequently, 142 original works are included in this review for an overview of these records in Section~\ref{sec:overview}, and their synthesis and discussion in Section~\ref{sec:discussion}. The information relative to the 12 data items for each of the included records is available in Appendix~\ref{a2:data-extraction} and also in the repository. The included works represent 34.55\% of the 411 eligible records for this review. This result indicates that the methodology followed in this review led to a high percentage of quality results.


\section{Results Overview}
\label{sec:overview}

In this section, the main goal is to overview the results not in terms of their scientific contribution but in terms of their bibliographic data for presenting an overview of the included records in the review. First, statistic results of the data sources in which the 144 included records could be identified in the methodology allow the evaluation of the coverage between the sources.
Next, the tool VOSviewer~\parencite{results:vosviewer:1,results:vosviewer:2} is used to obtain the co-occurrence analysis for the keywords and the authors. The former focus on the keywords recency and their occurrence in the sources, while the latter discusses the research networks between the authors, and the ones with more publications in long-term localization and mapping.
Lastly, two analysis are presented relative to the evolution of the publication year and most relevant publication venues.

\subsection{Data source}
\label{sec:overview:db}

The results on the identification phase are exported to BibTeX files from each data source. This exportation considers all the information available in the data sources, such as citation (e.g., author, title, publication venue, and type of record) and bibliographic (e.g., affiliation and the publisher) information of each record, the abstract, and author and indexed keywords. Next, using the \texttt{bibtexparser}\footnote{\url{https://bibtexparser.readthedocs.io/en/master/}} Python library, the BibTeX files are processed to identify uncompleted records. For example, the DOI must be specified and, if not available, the record's information must be manually completed with a corresponding URL. Then, considering the 144 included records in this review, a Python script searches each record in the BibTeX files corresponding to each data source. This search uses the DOI, URL, and title data to identify if a data source had in its identification results the searched record. Given that these three fields can contain lower and upper letters, the respective strings must be compared only after converting them to lower cases. As a result, the number of identified records by each data source of the 144 included ones in the review are the following ones:

\begin{itemize}[nosep]
\item \citetitle{methodology:search:db:acm}: 25 records (17.4\%);
\item \citetitle{methodology:search:db:dimensions}: 85 records (59.0\%);
\item \citetitle{methodology:search:db:ieee-xplore}: 68 records (47.2\%);
\item \citetitle{methodology:search:db:inspec}: 104 records (72.2\%);
\item \citetitle{methodology:search:db:scopus}: 122 records (84.7\%);
\item \citetitle{methodology:search:db:wos}: 105 records (72.9\%).
\end{itemize}

The database \citetitle{methodology:search:db:scopus} is the source that identified the greatest number of included records. This result was expected given that \citetitle{methodology:search:db:scopus} is considered as one of the largest curated databases~\parencite{methodology:search:db:coverage:dim-scopus-wos}, indexing more than 25000 active titles (e.g., conferences proceedings, journals) and 7000 publishers\footnote{\url{https://www.elsevier.com/solutions/scopus/how-scopus-works}}.
Two other sources with more than 70\% of identified records are \citetitle{methodology:search:db:inspec} and \citetitle{methodology:search:db:wos}. Similarly to \citetitle{methodology:search:db:scopus}, these two databases index also records from thousands of journals, conferences, and publishers\footnote{\url{https://www.elsevier.com/solutions/engineering-village/content/inspec}}\textsuperscript{,}\footnote{\url{https://clarivate.com/webofsciencegroup/solutions/web-of-science/}}.
Although \citetitle{methodology:search:db:dimensions} is also a bibliographic database covering millions of publications from thousands of sources, this database is the newest one (created in 2018) relative to the other three considered in this review (\citetitle{methodology:search:db:inspec}, \citetitle{methodology:search:db:scopus}, and \citetitle{methodology:search:db:wos}) and could be a factor to why it obtained a lower percentage (59.0\%) than the other three databases. Another possible reason is that \citetitle{methodology:search:db:scopus} and \citetitle{methodology:search:db:wos} have the majority of their coverage in Life Sciences, Physical Sciences, and Technology Area (including the Engineering subject area related to the topic of this review), while \citetitle{methodology:search:db:dimensions} has better coverage in Social Sciences and Arts \& Humanities~\parencite{methodology:search:db:coverage:dim-scopus-wos}.
Even though \citetitle{methodology:search:db:ieee-xplore} is a digital library and only indexes works published by IEEE and its partners, this data source returns 47.2\% of the include records in the review. The main reason is that this library indexes publications related to electrical engineering and computer science, subject areas related to long-term localization and mapping\footnote{\url{https://ieeexplore.ieee.org/Xplorehelp/overview-of-ieee-xplore/about-ieee-xplore}}.
Finally, the \citetitle{methodology:search:db:acm} using \textit{The ACM Guide to Computing Literature} collection only finds published records by ACM and possible links to other records focused exclusively on computing\footnote{\url{https://libraries.acm.org/digital-library/acm-guide-to-computing-literature}} and not directly related to the Computer Science or Engineering subject areas, explaining why this source obtained a lower coverage percentage of the included results than the other sources for this review.

Furthermore, Table~\ref{tab:overview:source} presents a coverage analysis of the identified results from each data source for the 144 included records in this review. Table~\ref{tab:overview:source:intersect} presents the pairwise overlap between sources. The corresponding percentage is the ratio of records identified by both sources to the one between the two that has the smallest number of results: $\#\{A\cap B\} / \text{min}\{\#A,\#B\}$, where $\#A$ and $\#B$ is the number of results for a data source $A$ and $B$, respectively, and $\#\{A\cap B\}$ is the intersection results between the two sources. For example, if the pairwise results is 100\%, it means that the data source with more records found was capable of obtaining all the results, i.e., had full coverage over the other source. Table~\ref{tab:overview:source:union} reports the percentage of records identified by at least one of two data sources over all 144 included records: $\#\{A\cup B\} / 144$, where $A\cup B$ is the union correspondence results of the sources $A$ and $B$. This percentage represents the joint coverage of two databases over the 144 included records.
					
\begin{table}[h]
  \centering
  \caption{Pairwise coverage analysis of the data sources considered in the review over the 144 included records: (a) identification only on both pairwise sources ($\#\{A\cap B\} / \text{min}\{\#A,\#B\}$); (b) on either ones ($\#\{A\cup B\} / \#\text{records}$). Legend: dim -- \citetitle{methodology:search:db:dimensions}, ieee -- \citetitle{methodology:search:db:ieee-xplore}, insp -- \citetitle{methodology:search:db:inspec}, scop -- \citetitle{methodology:search:db:scopus}, wos -- \citetitle{methodology:search:db:wos}.}
  \label{tab:overview:source}
  \subfloat[][]{%
  \label{tab:overview:source:intersect}%
  \begin{tabular}{l|P{0.08\columnwidth}P{0.08\columnwidth}P{0.08\columnwidth}P{0.08\columnwidth}P{0.08\columnwidth}P{0.08\columnwidth}}
\hline
$A\cap B$ & \textbf{acm} & \textbf{dim} & \textbf{ieee} & \textbf{insp} & \textbf{scop} & \textbf{wos}\\
\hline
\textbf{acm}  & -- & 96.0\% & 44.0\% & 88.0\% & 96.0\% & 96.0\%\\
\textbf{dim}  & -- & --     & 69.1\% & 77.6\% & 97.6\% & 95.3\%\\
\textbf{ieee} & -- & --     & --     & 89.7\% & 91.2\% & 73.5\%\\
\textbf{insp} & -- & --     & --     & --     & 87.5\% & 68.3\%\\
\textbf{scop} & -- & --     & --     & --     & --     & 89.5\%\\
\textbf{wos}  & -- & --     & --     & --     & --     & --\\
\hline
    \end{tabular}%
  }
  \linebreak
  \subfloat[][]{%
  \label{tab:overview:source:union}%
  \vspace{0.5em}
  \begin{tabular}{l|P{0.08\columnwidth}P{0.08\columnwidth}P{0.08\columnwidth}P{0.08\columnwidth}P{0.08\columnwidth}P{0.08\columnwidth}}
\hline
$A\cup B$ & \textbf{acm} & \textbf{dim} & \textbf{ieee} & \textbf{insp} & \textbf{scop} & \textbf{wos}\\
\hline
\textbf{acm}  & -- & 59.7\% & 56.9\% & 74.3\% & 85.4\% & 73.6\%\\
\textbf{dim}  & -- & --     & 73.6\% & 85.4\% & 86.1\% & 75.7\%\\
\textbf{ieee} & -- & --     & --     & 77.1\% & 88.9\% & 85.4\%\\
\textbf{insp} & -- & --     & --     & --     & 93.8\% & 95.8\%\\
\textbf{scop} & -- & --     & --     & --     & --     & 92.4\%\\
\textbf{wos}  & -- & --     & --     & --     & --     & -\\
\hline
  \end{tabular}%
  }
\end{table}

Analyzing the coverage results in Table~\ref{tab:overview:source}, the first observation is that the pairwise union results of two sources increase the independent coverage of each source. This observation validates the need identified in the methodology discussed in Section~\ref{sec:methodology} to consider several data sources in the identification phase of a review. Moreover, the pairwise union coverage of \citetitle{methodology:search:db:inspec}, \citetitle{methodology:search:db:scopus}, and \citetitle{methodology:search:db:wos} is greater than 90\% of the included records. When evaluating the joint coverage of these three databases, they identify all 144 of the included records, i.e., a 100\% coverage. Although this result could indicate that those three sources guarantee full coverage of the long-term localization and mapping research topic, it is always advisable to consider as most as possible sources in the methodology. Another observation is relative to the overlap of \citetitle{methodology:search:db:scopus} with the other sources, which is greater than 85\%. This overlap indicates that \citetitle{methodology:search:db:scopus} covers results not only on the topic of this review but also the results obtained by the other sources considered in the methodology. Lastly, \citetitle{methodology:search:db:inspec} and \citetitle{methodology:search:db:wos} achieve a pairwise overlap percentage of 68.3\% between themselves, while their union represents 95.8\% of the included records. This discrepancy indicates that these two sources identify unique results between themselves. Indeed, \citetitle{methodology:search:db:inspec} identifies 33/144 records not found by \citetitle{methodology:search:db:wos}, and vice-versa for \citetitle{methodology:search:db:wos}, with 34/144 unique records.

\subsection{Keywords co-occurrence}
\label{sec:overview:kw}

Next, VOSviewer~\parencite{results:vosviewer:1,results:vosviewer:2} is used to analyze the co-occurrence of keywords in the included articles. This co-occurrence is the relatedness of items determined based on the number of documents in which the keywords occur together.
For this analysis, first, a Python script processes the BibTeX file containing the citation and bibliographic information, the author and the indexed keywords, and the abstract of the records to join the author with the indexed keywords in the same \texttt{keywords} field. Then, an online tool\footnote{\url{https://www.bibtex.com/c/bibtex-to-ris-converter/}} converts this processed BibTeX to a RIS file.
Even though VOSviewer supports file types directly exported from \citetitle{methodology:search:db:dimensions}, \citetitle{methodology:search:db:scopus}, or \citetitle{methodology:search:db:wos} as input, none of these data sources obtained all the 144 included records of the review in the identification phase. Given that VOSviewer does not support BibTeX files, the conversion to RIS file is required for using as input.
The disadvantage of using this file format in VOSviewer is only allowing to perform co-occurrence of items (e.g., keywords or authors), while bibliographic data from \citetitle{methodology:search:db:dimensions}, \citetitle{methodology:search:db:scopus}, or \citetitle{methodology:search:db:wos} in CSV files would allow other analysis such as citation, co-citation, or bibliographic coupling.
However, the creation of these CSV files follow different templates depending on the data source. So, RIS files allow the integration of all 144 included records for obtaining the two co-occurrence analysis presented in this review (namely, keywords and co-authorship).

In Figure~\ref{fig:overview:kw:original}, the network presents the overlay visualization of the keywords co-occurrence in the included records weighted by the number of occurrences of each term, using full counting for the links' strength. The latter computes the strength of the links directly by the number of co-occurrences of the respective two terms.
The overlay visualization colors the keywords differently according to the average publication year of the included records in which each of the keywords appears. This coloring allows analyzing which are the ones that are associated with the most recent publications.
As for the keywords' weighting, the number of occurrences dictates the size of the circles.
Furthermore, the minimum number of occurrences of a keywords set in VOSviewer for obtaining the graph is 5 originating the 35 keywords illustrated in Figure~\ref{fig:overview:kw:original}. This parameter was selected for visualization purposes while also filtering uninteresting keywords.
Similarly, setting the attraction and repulsion parameters to 2 and 0, respectively, distances the terms more from each other than using the values recommended in the VOSviewer manual\footnote{\url{https://www.vosviewer.com/documentation/Manual_VOSviewer_1.6.8.pdf}} (2 and 1, respectively). These two parameters only interfere in the localization of the terms in the map, not in the graph connections.
Lastly, a thesaurus of the keywords (available in the repository) is used to join similar terms: spelling differences (e.g., localization -- localisation), full terms versus abbreviations (simultaneous localization and mapping -- SLAM), while also allowing the concatenation of long keywords for visualization reasons.

\begin{figure*}[!t]
  \centering
  \subfloat[][]{\includegraphics[width=\columnwidth]{figures/kw.png}%
  \label{fig:overview:kw:original}}
  \hfill
  \subfloat[][]{\includegraphics[width=\columnwidth]{figures/kw_long-term.png}%
  \label{fig:overview:kw:long-term}}
  \caption{Keywords co-occurrence analysis on the 144 included records generated by VOSviewer with overlay visualization by the average publication year: (a) original keywords; (b) all keywords containing long-term and lifelong summarized by the terms themselves. Parameters used for generating the co-occurrence network: minimum number of occurrences = 5, attraction = 2, repulsion = 0, scale = 1.49, circles size variation = 0.5, lines size validation = 1.0. Legend: \texttt{cnn} -- Convolutional Neural Networks, \texttt{env} -- environment, \texttt{localiz} -- localization, \texttt{nav} -- navigation, \texttt{nn} -- Neural Networks, \texttt{recog} -- recognition, \texttt{vslam} -- visual SLAM.}
  \label{fig:overview:kw}
\end{figure*}

Overall, the keyword \texttt{robot} is the one that appears more times in the included records: 111 occurrences, links with 34 other terms, and has a total link strength of 403 (sum of co-occurrences of all of its links). This result is expected due to the relation of this review's topic to robotics.
Similarly, three other keywords in the network related to long-term localization and mapping topic with high values of occurrence, number of links, and total link strength are \texttt{slam} (75, 34, and 288), \texttt{mapping} (48, 33, and 204), and \texttt{localization} (47, 32, and 194, respectively). The methodology for the search strategy discussed in Section~\ref{sec:methodology:search} considers all of these four keywords. Thus, the significant influence of \texttt{robot}, \texttt{slam}, \texttt{mapping}, and \texttt{localization} in the keywords co-occurrence analysis indicates that, after the all the phases executed in this review's methodology, the 144 included records have a high correlation with the keywords considered in the search query. Given that the keywords are usually selected or indexed to capture the essence of the document, this correlation indicates that the search query is appropriate to obtain the search results, even considering only the keywords as search fields.

As for keywords related to the outcome of the PICO framework, \texttt{long-term autonomy} occurs only 6 times in the included records, linking with 16 other keywords and having a total link strength of 27. This low occurrence could indicate that the term \texttt{long-term autonomy} is not usually used by the authors nor indexed by the databases. 
However, the specific term of \texttt{long-term autonomy} does not summarize all the possibilities for the outcome of the PICO framework (see Section~\ref{sec:purpose}). Indeed, for this reason, the search query for the identification phase uses only the following single terms: \texttt{"long term"} and \texttt{"life long"} (resumes the possibility of having a space or a hyphen), and \texttt{lifelong}.
Figure~\ref{fig:overview:kw:long-term} presents the keywords co-occurrence analysis using the same parameters for obtaining Figure~\ref{fig:overview:kw:original}. The difference to the latter network is using a thesaurus that summarizes all the keywords that contain \texttt{long-term} and \texttt{lifelong} into the terms themselves, obtaining 36 keywords with a minimum of 5 occurrences in the 144 included records.
In terms of occurrences, number of links, and total link strength, the impact of the thesaurus keyword \texttt{long-term} is 25, 28, and 105, and for lifelong 6, 17, and 31, respectively. These values are much higher than the ones respective only to \texttt{long-term autonomy} from Figure~\ref{fig:overview:kw:original}.
The reason is that \texttt{long-term} in Figure~\ref{fig:overview:kw:long-term} compiles the occurrences of keywords such as \texttt{long-term autonomy}, \texttt{long-term mapping}, and \texttt{long-term localization} (6, 2, and 2 occurrences, respectively), and \texttt{lifelong} sum up, for example, three different versions of \texttt{lifelong learning} (using \texttt{lifelong}, \texttt{life-long} and \texttt{life long} with 2, 1, and 2 occurrences, respectively) and \texttt{lifelong slam} (1 occurrence). Hence, these results proves that the third \texttt{AND} part of the search query (\texttt{"long term" OR "life long" OR lifelong}) covers well the PICO framework's outcome. Plus, they also show no consensus among the authors and by the databases indexation on how to define a keyword for the topic of long-term localization and mapping.

In terms of the average year of publication, analyzing the diagrams in Figure~\ref{fig:overview:kw} on its colorization, the first observation is the recency of terms related to visual localization. The keywords visual SLAM (\texttt{vslam}), visual navigation (\texttt{visual nav}), and visual localization (\texttt{visual localiz}) have all an average publication year higher than 2017. This recency indicates that recent approaches related to the topic of this review, long-term localization and mapping, are more inclined to use vision as a sensorization input.
Another sensor that appeared with high relevance in the network is \texttt{radar}, with 15 occurrences and an average publication year of 2019.20. This sensor is agnostic to the environment changes such as illumination and season changes intrinsically associated with vision and could be the reason why the recent works related to long-term localization and mapping are using it.
Moreover, place recognition (\texttt{place recog}) stands out not only by its recency but importance. The keyword itself (\texttt{place recog}) occurs 31 times and an average publication year of 2017.77, with terms related to place recognition such as \texttt{loop closure} and global localization (\texttt{global localiz}) with recent average publication years (2017.82 and 2018.75, respectively) and strong link to place recognition (5 co-occurrences for each of the links between \texttt{loop closure} and \texttt{global localiz} with \texttt{place recog}). Lastly, machine learning also seems to be used in recent works included in this review. The keywork learning occurs 7 times with an average publication year of 2017.00. Neural Netowrks (\texttt{nn}), Convolutional Neural Networks (\texttt{cnn}), and \texttt{deep learning} have a similar number of occurrences (6, 5, and 8) and publication years higher than 2017 (2017.83, 2018.00, and 2019.12, respectively). These results could mean another trend of using machine learning to improve the long-term autonomy of mobile robots.

Although the recency of keywords related to dynamic environments is lower than 2017 (2015.50 and 2016.75 for \texttt{dynamic env} and \texttt{env change}), they have a high occurrence (14 and 12, respectively), located close to each other in the network, and have a strong link between them (4 co-occurrences). Three keywords also located near each other are \texttt{graph}, \texttt{marginalization}, and information theory (\texttt{info theory}) while having similar average publication years (2014.36, 2014.00, and 2015.50, respectively). Even though the number of occurrences of these terms is low (11, 6, and 5 for \texttt{graph},\texttt{ marginalization}, and \texttt{info theory}, respectively), their map proximity could indicate a focus in the past on the topic of graph sparsity, i.e., maintaining the graph in the long-term to only depend on the environment size and not on the robot's operation time.

The keywords co-occurrence analysis also relates to the categories of DE1 (see Section~\ref{sec:methodology:data}). Works associated with place recognition, global localization, and loop closure terms require invariance to the appearance changes in the environment, equivalent to the appearance category. The dynamics category is associated with works focused on dynamic environments. As for the other group of keywords with a high occurrence and strong links between each other, the ones related to graph and information theory, the respective works focus on removing uninformative data from the map~\parencite{kretzschmar-stachniss:2012:0278364912455072}, which is related to map sparsification, and so, to the sparsity category of DE1. These relations between the appearance, dynamics, and sparsity categories to the semantic analysis of the keywords co-occurrence supports the categorization of DE1 considered in this review, while also indicating that the discussion on the proposed methodologies should focus on each one of the categories. Even though the two remaining categories of DE1 (multi-session and computational) are not represented in the keyword analysis, the execution of the data extraction phase identified the need for having these two categories, given the importance of multi-session handling and computational efficiency for long-term localization and mapping. However, each category of DE1 will be discussed in Section~\ref{sec:discussion} in further detail.


\subsection{Co-authorship analysis}
\label{sec:overview:authors}

The other analysis obtained using VOSviewer is the co-authorship network presented in Figure~\ref{fig:overview:authors}. Similar to the keywords network illustrated in Figure~\ref{fig:overview:kw}, the co-occurrence of the authors' names creates links among them in the graph. The strength of these links is dictated by the number of documents the two authors of a link are co-authors in the same record, and the number of co-authored works determines the size of the circles respective to each author in the graph. 
In contrast to Figure~\ref{fig:overview:kw}, the network in Figure~\ref{fig:overview:authors} does not have any overlay specific to coloring depending on the average publication year. Instead, the main goal of the co-authorship analysis in this review is to present possible research networks detected in the 144 included records. Thus, the coloring in Figure~\ref{fig:overview:authors} represents the clusters of authors detected by VOSviewer. This network only considers authors with a minimum of 3 works for relevance and visualization reasons, resulting in 29 authors.
Also, authors identified only by the initial of the first name and by the surname can lead to incorrect correspondences in terms of co-authorship. VOSviewer detects 392 authors in the 144 included records using the original RIS file used in Section~\ref{sec:overview:kw} compared to 413 after checking the authors names. Indeed, a manual check is performed on all authors of the included records to guarantee no false correspondences for the co-authorship analysis with VOSviewer. This manual check ensures each author has its full first and surname and any middle initials while also using the same name for an author in different records.

\begin{figure}[h]
  \centering
  \includegraphics[width=\columnwidth]{figures/authors.png}
  \caption{Co-authorship analysis on the 144 included records generated by VOSviewer. Parameters used for generating the co-occurrence network: minimum number of occurrences = 3, attraction = 2, repulsion = -3, scale = 1.49, circles size variation = 1.0, lines size validation = 1.0.}
  \label{fig:overview:authors}
\end{figure}

Analyzing Figure~\ref{fig:overview:authors}, the co-authorship network presents 8 clusters. These clusters are separated from each other, i.e., no link exists between authors from different clusters. However, this separation does not mean that there is not any co-authorship between authors from different clusters only indicating that for a miminum of 3 co-authored documents there is not a connection between these 8 clusters. Even so, the graph presented in Figure~\ref{fig:overview:authors} allows the identification of the most relevant research networks in terms of number of co-authored documents and in the context of long-term localization and mapping, considering the 144 records included in this review. As a results, the following enumeration presents the authors that belong to each cluster in the format of author (number of co-authored documents):

\begin{enumerate}\setlength\itemsep{-0.5em}
\item Rong Xiong \orcid{0000-0001-9318-9014}\textsuperscript{,}\scholar{1hI9bqUAAAAJ} (7),
        Yue Wang \orcid{0000-0002-0981-935X}\textsuperscript{,}\scholar{N543LSoAAAAJ} (7),
        Huan Yin \orcid{0000-0002-0872-8202}\textsuperscript{,}\scholar{1fNc3vUAAAAJ} (6),
        Li Tang \orcid{0000-0003-2590-6872} (5),
        and Xiaqing Ding \orcid{0000-0001-7802-0130}\textsuperscript{,}\scholar{6u5OHUcAAAAJ} (4);
\item Hao Zhang \scholar{Ug2VxyUAAAAJ} (6),
        John J. Leonard \orcid{0000-0002-8863-6550}\textsuperscript{,}\scholar{WPe7vWwAAAAJ} (5),
        Michael Kaess \scholar{27eupmsAAAAJ} (5),
        Fei Han \orcid{0000-0002-8619-3987} (3),
        Guoquan Huang \scholar{trMUyZIAAAAJ} (3),
        José Neira \orcid{0000-0003-0668-977X}\textsuperscript{,}\scholar{scoMbR8AAAAJ} (3),
        and Ryan M. Eustice \orcid{0000-0002-9989-4942}\textsuperscript{,}\scholar{WroYmiAAAAAJ} (3);
\item Wolfram Burgard \orcid{0000-0002-5680-6500}\textsuperscript{,}\scholar{zj6FavAAAAAJ} (7),
        Cyrill Stachniss \orcid{0000-0003-1173-6972}\textsuperscript{,}\scholar{8vib2lAAAAAJ} (5),
        Giorgio Grisetti \orcid{0000-0002-8038-9989}\textsuperscript{,}\scholar{yD-SFG4AAAAJ} (3),
        Henrik Kretzschmar \scholar{pZuaUlwAAAAJ} (3),
        and Tayyab Naseer \orcid{0000-0002-3350-3005}\textsuperscript{,}\scholar{1FePZqEAAAAJ} (3);
\item Roland Siegwart \orcid{0000-0002-2760-7983}\textsuperscript{,}\scholar{MDIyLnwAAAAJ} (6),
        Igor Gilitschenski \orcid{0000-0001-6426-365X}\textsuperscript{,}\scholar{Nuw1Y4oAAAAJ} (3),
        Marcin Dymczyk \orcid{0000-0003-3667-8764}\textsuperscript{,}\scholar{XYHy7U8AAAAJ} (3),
        Michael Bosse \scholar{eopb1VgAAAAJ} (3),
        and Paul Furgale \orcid{0000-0002-7367-1046}\textsuperscript{,}\scholar{RNDtSG8AAAAJ} (3);
\item Tom Duckett \orcid{0000-0003-2971-7905}\textsuperscript{,}\scholar{et1GU2EAAAAJ} (6),
        Cheng Zhao \orcid{0000-0001-8502-3233}\textsuperscript{,}\scholar{EAC-8m0AAAAJ} (3),
        and Tomáš Krajník \orcid{0000-0002-4408-7916}\textsuperscript{,}\scholar{Qv3nqgsAAAAJ} (3);
\item Gordon F. Wyeth \orcid{0000-0002-4996-3612}\textsuperscript{,}\scholar{yfXZfXEAAAAJ} (5)
        and Michael J. Milford \orcid{0000-0002-5162-1793}\textsuperscript{,}\scholar{TDSmCKgAAAAJ} (4);
\item Paul Newman \scholar{BtO5fTUAAAAJ} (4);
\item Timothy D. Barfoot \orcid{0000-0003-3899-631X}\textsuperscript{,}\scholar{N_vPIhoAAAAJ} (3).
\end{enumerate}

When analyzing the affiliations of the authors mentioned previously at the time of publication, all authors of the first cluster belonged to the State Key Laboratory of Industrial Control and Technology (SKLICT) and the Institute of Cyber-Systems and Control at Zhejiang University in China. Even though Huan Yin, Yue Wang, Xiaqing Ding, Li Tang, and Rong Xiong mention their affiliation to the Joint Centre for Robotics Research between Zhejiang University, China, and the University of Technology Sydney, Sydney, in the work \parencite{yin-et-al:2020:2905046}, this specific affiliation only appeared in this article. The total link strength (sum of all links weights) of each of the authors in that cluster is higher than 16, meaning a high co-authorship between them. Indeed, all five authors have links between all of them.
Similar to the first cluster, the third, fourth, fifth, and sixth clusters have common affiliations within each one: the Autonomous Intelligent Systems at the University of Freiburg in Germany, the Autonomous Systems Lab (ASL) at ETH Zürich in Switzerland, the Lincoln Centre for Autonomous Systems (LCAS) at the University of Lincoln in UK, and the School of Electrical Engineering and Computer Science at Queensland University of Technology (QUT) in Australia.
However, the interlinking between the authors is not as strong as in the first cluster, as shown in Figure~\ref{fig:overview:authors} by the authors of these clusters not being connected between all the ones within each cluster. Even so, the common affiliation shows there is considerable interest by these research units in the long-term localization and mapping topic.

The affiliation analysis in the second cluster is more complex given that there was no affiliation common to all authors at the time of the records' publication. Instead, the following affiliations were found: Fei Han and Hao Zhang with the Department of Computer Science at Colorado School of Mines in the USA, Guoquan Huang with the Department of Mechanical Engineering at the University of Delaware in the USA, John J. Leonard and Michael Kaess with the Computer Science and Artificial Intelligence Laboratory (CSAIL) at the Massachusetts Institute of Technology (MIT) in the USA, Ryan M. Eustice with the Perceptual Robotics Laboratory (PeRL) at the University of Michigan in the USA, and José Neira with the Instituto Universitario de Investigación en Ingeniería de Aragón (I3A) at the Universidad de Zaragoza in Spain.
Although there are 5 different affiliations to which the 7 authors stated in the respective records, 4 of the research institutions noted for the second cluster are in the USA, indicating a possible reason for facilitating the linkage between these authors from different research units.

In terms of the clusters composed by single authors, the affiliations of Paul Newman and Timothy D. Barfoot are the Oxford Robotics Institute at the University of Oxford in UK and the Autonomous Space Robotics Laboratory (ASRL) at the University of Toronto Institute for Aerospace Studies (UTIAS) in Canada, respectively. Even though these two authors are not linked with any others in the network, the co-authorship analysis indicates that they have an interest in long-term localization and mapping. This interest is shown by their number of co-authored records: 4 and 3 by Paul Newman and Timothy D. Barfoot, respectively.

As for the number of co-authored publications, considering the 144 included records, the authors that appeared to have more research on the review's topic are Rong Xiong, Yue Wang, and Wolfram Burgard, given the 7 co-authored publications of each one. However, Rong Xiong and Yue Wang have co-authored the 7 documents attributed to each of them. This relation and similar ones can biase the analysis of which authors are having more impact in the review's topic.
The clustering shown in Figure~\ref{fig:overview:authors} allows a more unbiased analysis relative to the co-authorship links between authors. Thus, based on the clustering and which author from each cluster has the most co-authored publications, the most influential authors in long-term localization and mapping are the following ones: Rong Xiong (or Yue Wang), Hao Zhang, Wolfram Burgard, Roland Siegwart, Tom Duckett, Gordon F. Wyeth, Paul Newman, and Timothy D. Barfoot.

\subsection{Year of publication}
\label{sec:overview:year}

The relevance of the long-term localization and mapping topic can be evaluated by the evolution of the number of publications. Figure~\ref{fig:overview:year} presents this evolution from the earliest year of publication of the included records to the year at the time of writing this article. The latter has its respective data dashed to indicate that the last year is not completed at the time of writing. Analyzing Figure~\ref{fig:overview:year}, this review's topic seems to have gain relevance in 2009 with 6 works, compared to only one publication in 2007 and another in 2002 in the previous years to 2009. From that year onwards, the graph has an almost linear tendency reaching a maximum of 23 records in 2021, while already having 8 publications in 2022 until May 17, 2022. This tendency shows that long-term localization and mapping is gaining interest throughout the years and, consequently, supports the importance and relevance of this review for the scientific community.

\begin{figure}[!t]
  \centering
  \includegraphics[width=\columnwidth]{figures/year.png}
  \caption{Evolution of published records per year considering the 144 included records in this review}
  \label{fig:overview:year}
\end{figure}

\subsection{Publication venue}
\label{sec:overview:publication}

Finally, the last overview of the 144 included records in the review is relative to the publication venue. Table~\ref{tab:overview:publication} presents the venues with more than 1 publication, separating the journals and conferences in two different tables (Tables~\ref{tab:overview:publication:journal} and \ref{tab:overview:publication:conference}, respectively).
The columns $\mu$ present the average year of publication of the records associated to a certain venue, while $\text{max}$ columns display the publishing recency by the year of the most recent publication in the venue. For comparing to the average value ($\mu$), the third column ($\sigma$) of each table presents the standard deviation based on the publication year data.
The last column state the number of records published in the venue from the 144 records included in the review for discussion.

\begin{table}[!h]
  \centering
  \caption{Publication venues of the included records in this review with more than one record published in the venue: (a) journals; (b) conferences. Legend: $\mu$ -- average year of publication, $\sigma$ -- standard deviation of the publication year, $\text{max}$ -- maximum year of publication, \# -- number of records published at a certain venue}
  \label{tab:overview:publication}
  \subfloat[][]{%
  \begin{tabular}{p{0.55\columnwidth}cccc}
\hline
                 & \multicolumn{3}{c}{\textbf{Year}} & \\
\cline{2-4}
\textbf{Journal} & $\mu$ & $\sigma$ & $\text{max}$ & \textbf{\#}\\
\hline
Robotics and Autonomous Systems & 2016 & 3.9 & 2021 & 13\\
IEEE Robotics and Automation Letters & 2019 & 1.7 & 2022 & 12\\
International Journal of Robotics Research & 2014 & 3.2 & 2022 & 11\\
Journal of Field Robotics & 2017 & 3.5 & 2022 & 8\\
Autonomous Robots & 2017 & 2.2 & 2020 & 7\\
IEEE Transactions on Intelligent Transportation Systems & 2021 & 0.8 & 2022 & 4\\
Sensors & 2019 & 0.8 & 2020 & 4\\
IEEE Transactions on Robotics & 2017 & 3.1 & 2022 & 4\\
IEEE Sensors Journal & 2020 & 1.5 & 2021 & 2\\
International Journal of Advanced Robotic Systems & 2020 & 1.5 & 2021 & 2\\
\hline
  \end{tabular}\label{tab:overview:publication:journal}%
  }
  \linebreak
  \subfloat[][]{%
  \begin{tabular}{p{0.55\columnwidth}cccc}
\hline
                    & \multicolumn{3}{c}{\textbf{Year}} & \\
\cline{2-4}
\textbf{Conference} & $\mu$ & $\sigma$ & $\text{max}$ & \textbf{\#}\\
\hline
IEEE International Conference on Robotics and Automation (ICRA) & 2016 & 3.9 & 2021 & 22\\
IEEE/RSJ International Conference on Intelligent Robots and Systems (IROS) & 2017 & 3.8 & 2021 & 18\\
IEEE International Conference on Robotics and Biomimetics (ROBIO) & 2019 & 2.1 & 2021 & 3\\
IEEE International Intelligent Transportation Systems Conference (ITSC) & 2018 & 2.4 & 2021 & 3\\
European Conference on Mobile Robots (ECMR) & 2014 & 0.9 & 2015 & 3\\
IEEE Intelligent Vehicles Symposium (IV) & 2019 & 0.5 & 2019 & 2\\
International Conference on 3D Vision (3DV) & 2018 & 1.5 & 2019 & 2\\
International Conference on Advanced Robotics (ICAR) & 2011 & 2.0 & 2013 & 2\\
\hline
  \end{tabular}\label{tab:overview:publication:conference}%
  }
\end{table}

In terms of journals, the Robotics and Autonomous Systems, IEEE Robotics and Automation Letters, and the International Journal of Robotics stand out with more than 10 publications. Also, these journals have a high standard deviation (greater than 1.5), indicating that the publications spread out throughout the years.
In the case of the IEEE Robotics and Automation Letters, these results gain more relevance indicating a recent trend on publishing on this journal, considering that its creation was only on 2015\footnote{\url{https://www.ieee-ras.org/publications/ra-l}}.
With more than 5 publications, the Journal of Field Robotics and the Autonomous Robots have recent average of publication (2017) with a high standard deviation (greater than 2.0), similarly indicating that authors have been publishing in these two journals along the years.
In contrast, the IEEE Transactions on Intelligent Transportation Systems and Sensors journals have a standard deviation lower than 1 year, with an average publication year of at least 2019. The recency of publication on these two journals with a very low deviation suggests a recent interest of the authors to publish in these two journals works related to long-term localization and mapping.

As for conferences, the data in Table~\ref{tab:overview:publication:conference} shows a high discrepancy in the number of publications related to this review's topic in ICRA and IROS compared to the other venues. Indeed, all the other conferences have only a maximum of 3 records published in them, compared to 22 and 18 papers in ICRA and IROS, respectively. When considering that 62 of the 144 included records are published in conferences, ICRA and IROS with a total of 40 published works related to this review's topic represent 65\% of works published in conferences and 27.8\% of all included records. This result expresses the high relevance of ICRA and IROS in the topic of long-term localization and mapping.


\section{Discussion}
\label{sec:discussion}

\subsection{Appearance variance}
\label{sec:discussion:appearance}

\subsection{Dynamic environments}
\label{sec:discussion:dynamics}

\subsection{Map sparsification}
\label{sec:discussion:sparsify}

\subsection{Multi-session}
\label{sec:discussion:multisession}

\subsection{Computational}
\label{sec:discussion:computational}

\subsection{Evaluation metrics}
\label{sec:discussion:metrics}

\subsection{Long-term experimental data}
\label{sec:discussion:experiments}

\onecolumn

\renewcommand{\arraystretch}{1.125}
\setlength{\tabcolsep}{1.5pt}

\begin{tiny}

\begin{longtable}{@{\extracolsep{1pt}}
p{0.065\textwidth}                          % Dataset name
c|                                          % Year
cccccc|                                     % Long-term
p{0.095\textwidth}                          % Environment
p{0.095\textwidth}|                         % Domain
c|                                          % Sensory setup: odo
ccccccc|                                    %                camera
cc|                                         %                laser
cccc|                                       %                radar, sonar, imu, gps
cc|                                         % Sensor calibrations
p{0.095\textwidth}p{0.095\textwidth}|       % Ground-truth + Format
ccccc@{}}                                   % Distance, time, \#sequences
  \caption{Datasets used in the 144 included records for long-term localization and / or mapping experiments. Legend: odo -- odometry (wheeled, laser, visual, inertial, or a combination of odometry sources, dist. -- total distance length of the dataset, path -- total path distance if repeated several times, time -- total operation time, int. -- time interval between the start and end acquisition dates / time instants (d/w/m/y equivalent to day/week/month/year, 0 if only 1 run), and seq. -- number of sequences of the dataset.}
  %\vspace{0.5em}
  \label{tab:discussion:datasets}\\

%% FIRST TABLE HEADER
\hline
&&
\multicolumn{6}{c|}{\textbf{Long-term}} &
&&
\multicolumn{14}{c|}{\textbf{Sensor}} &
\multicolumn{2}{c|}{\textbf{Calib.}} &
&&
&&&&\\
&&&&&&&&&&
\multicolumn{1}{c}{} &
\multicolumn{7}{c}{\textbf{camera}} &
\multicolumn{2}{c}{\textbf{laser}}  &
&&&&
&&
&&
&&&&\\
\cline{12-18}
\cline{19-20}
\textbf{Dataset} & \textbf{Year} &
\rotatebox{90}{\textbf{lighting}} & \rotatebox{90}{\textbf{day/night}} & \rotatebox{90}{\textbf{weather}} & \rotatebox{90}{\textbf{seasonal}} & \rotatebox{90}{\textbf{dynamics}} & \rotatebox{90}{\textbf{sparsity}} & 
\textbf{Environ.} & \textbf{Domain} & 
\multicolumn{1}{c}{\rotatebox{90}{\textbf{odo}}} &
\rotatebox{90}{\textbf{gray}} & \rotatebox{90}{\textbf{color}} & \rotatebox{90}{\textbf{monocular}} & \rotatebox{90}{\textbf{stereo}} & \rotatebox{90}{\textbf{omni}} & \rotatebox{90}{\textbf{RGBD}} & \multicolumn{1}{c}{\rotatebox{90}{\textbf{thermal}}} &
\rotatebox{90}{\textbf{2D}} & \multicolumn{1}{c}{\rotatebox{90}{\textbf{3D}}} &
\rotatebox{90}{\textbf{radar}} & \rotatebox{90}{\textbf{sonar}} & \rotatebox{90}{\textbf{IMU}} & \rotatebox{90}{\textbf{GPS}} &
\rotatebox{90}{\textbf{intrinsic}} & \rotatebox{90}{\textbf{extrinsic}} & 
\textbf{GT data} & \textbf{Format} & 
\rotatebox{90}{\textbf{dist.} (km)} & \rotatebox{90}{\textbf{path} (km)} & \rotatebox{90}{\textbf{time} (h)} & \rotatebox{90}{\textbf{int.} (d/w/m/y)} & \rotatebox{90}{\textbf{\#seq.}}\\
\hline
\endfirsthead

%% TABLE HEADER IN THE FOLLOWING PAGES
\multicolumn{27}{l}{\itshape{\tablename\ \thetable{}: continued from previous page}}
\vspace{0.5em}\\
\hline
&&
\multicolumn{6}{c|}{\textbf{Long-term}} &
&&
\multicolumn{14}{c|}{\textbf{Sensor}} &
\multicolumn{2}{c|}{\textbf{Calib.}} &
&&
&&&&\\
&&&&&&&&&&
\multicolumn{1}{c}{} &
\multicolumn{7}{c}{\textbf{camera}} &
\multicolumn{2}{c}{\textbf{laser}}  &
&&&&
&&
&&
&&&&\\
\cline{12-18}
\cline{19-20}
\textbf{Dataset} & \textbf{Year} &
\rotatebox{90}{\textbf{lighting}} & \rotatebox{90}{\textbf{day/night}} & \rotatebox{90}{\textbf{weather}} & \rotatebox{90}{\textbf{seasonal}} & \rotatebox{90}{\textbf{dynamics}} & \rotatebox{90}{\textbf{sparsity}} & 
\textbf{Environ.} & \textbf{Domain} & 
\multicolumn{1}{c}{\rotatebox{90}{\textbf{odo}}} &
\rotatebox{90}{\textbf{gray}} & \rotatebox{90}{\textbf{color}} & \rotatebox{90}{\textbf{monocular}} & \rotatebox{90}{\textbf{stereo}} & \rotatebox{90}{\textbf{omni}} & \rotatebox{90}{\textbf{RGBD}} & \multicolumn{1}{c}{\rotatebox{90}{\textbf{thermal}}} &
\rotatebox{90}{\textbf{2D}} & \multicolumn{1}{c}{\rotatebox{90}{\textbf{3D}}} &
\rotatebox{90}{\textbf{radar}} & \rotatebox{90}{\textbf{sonar}} & \rotatebox{90}{\textbf{IMU}} & \rotatebox{90}{\textbf{GPS}} &
\rotatebox{90}{\textbf{intrinsic}} & \rotatebox{90}{\textbf{extrinsic}} & 
\textbf{GT data} & \textbf{Format} & 
\rotatebox{90}{\textbf{dist.} (km)} & \rotatebox{90}{\textbf{path} (km)} & \rotatebox{90}{\textbf{time} (h)} & \rotatebox{90}{\textbf{int.} (d/w/m/y)} & \rotatebox{90}{\textbf{\#seq.}}\\
\hline
\endhead

%% TABLE FOOTER
\hline
\endfoot
\hline
\endlastfoot



\citetitle{dataset:fhw} & 2001 &  &  &  &  &  & x & indoor (museum) & ground (TOURBOT) & x &  &  &  &  &  &  &  & x &  &  &  &  &  &  &  & -- & CARMEN & -- & -- & 1.98 & -- & 1\\
\hline
\citetitle{dataset:fr079} & 2003 &  &  &  &  &  & x & indoor (office) & ground (robot) & x &  &  &  &  &  &  &  & x &  &  &  &  &  &  &  & -- & CARMEN & -- & -- & 0.29 & -- & 1\\
\hline
\citetitle{dataset:fr101} & 2003 &  &  &  &  &  & x & indoor (office) & ground (robot) & x &  &  &  &  &  &  &  & x &  &  &  &  &  &  &  & -- & CARMEN & -- & -- & 0.29 & -- & 1\\
\hline
\citetitle{dataset:intel03} & 2003 &  &  &  &  &  & x & indoor (office) & ground (robot) & x &  &  &  &  &  &  &  & x &  &  &  &  &  &  &  & -- & CARMEN & 0.506 & -- & 0.75 & -- & 1\\
\hline
\citetitle{dataset:mit-kilian} & 2004 &  &  &  &  &  & x & indoor (office) & ground (robot) & x &  &  &  &  &  &  &  & x &  &  & x &  &  &  &  & -- & CARMEN & 2.2 & -- & 2.5 & -- & 1\\
\hline
\citetitle{dataset:city-center-fabmap} & 2008 & x &  &  &  & x &  & outdoor (urban) & ground (robot) &  &  & x & x &  &  &  &  &  &  &  &  &  & x & x &  & GPS, manual & plain text (non-img), jpg (img) & 2 & -- & -- & -- & 1\\
\hline
\citetitle{dataset:lip6indoor} & 2008 &  &  &  &  &  &  & indoor (office) & ground (handheld) &  &  & x & x &  &  &  &  &  &  &  &  &  &  & x &  & manual & ppm (imgs) & -- & -- & 0.11 & -- & 1\\
\hline
\citetitle{dataset:lip6outdoor} & 2008 & x &  &  &  & x &  & outdoor (campus) & ground (handheld) &  &  & x & x &  &  &  &  &  &  &  &  &  &  & x &  & manual & ppm (imgs) & -- & -- & 0.3 & -- & 1\\
\hline
\citetitle{dataset:new-college-fabmap} & 2008 & x &  &  &  & x &  & outdoor (campus) & ground (robot) &  &  & x & x &  &  &  &  &  &  &  &  &  & x & x &  & GPS, manual & plain text (non-img), jpg (img) & 1.9 & -- & -- & -- & 1\\
\hline
\citetitle{dataset:stlucia07} & 2008 & x &  &  &  & x &  & outdoor (urban) & ground (car) &  & x &  & x &  &  &  &  &  &  &  &  &  &  &  &  & -- & -- & 66 & -- & 1.67 & -- & 1\\
\hline
\citetitle{dataset:bicocca-indoor} & 2009 & x &  &  &  & x & x & indoor (office) & ground (Robocom) & x & x & x & x & x & x &  &  & x &  &  & x & x &  & x & x & map model, laser-based & plain text (non-img), png (img) & -- & -- & 2.5 & 3d & 5\\
\hline
\citetitle{dataset:cold} & 2009 & x & x & x &  & x &  & indoor (office) & ground (Pioneer 3, ATRV Mini, PeopleBot) & x &  & x & x &  & x &  &  & x &  &  &  &  &  &  &  & laser-based, manual & plain text (non-img), jpg (img) & 0.92 & -- & 0.99 & -- & 76\\
\hline
\citetitle{dataset:malaga-2009} & 2009 & x &  &  &  & x &  & outdoor (parking, campus) & ground (car) &  &  & x & x &  &  &  &  & x &  &  &  & x & x & x & x & RTK-GPS & Rawlog MRPT & 6.358 & -- & -- & -- & 6\\
\hline
\citetitle{dataset:newcollege} & 2009 & x &  &  &  & x &  & outdoor (campus) & ground (Segway) & x & x & x &  & x & x &  &  & x &  &  &  & x & x & x & x & GPS & plain text (non-img), png, jpg (img) & 2.2 & -- & 0.73 & -- & 1\\
\hline
\citetitle{dataset:albert-b} & 2010 & x &  &  &  &  & x & indoor (office) & ground (iRobot B21r) & x &  & x & x &  &  &  &  & x &  &  &  &  &  &  &  & -- & CARMEN (non-img), jpg (img) & -- & -- & 0.18 & -- & 1\\
\hline
\citetitle{dataset:cmu-vl} & 2011 & x &  &  & x & x &  & outdoor (urban) & ground (car) &  &  & x & x &  &  &  &  &  &  &  &  &  & x & x & x & GPS & -- & -- & 8.5 & -- & 1y & 16\\
\hline
\citetitle{dataset:ford} & 2011 & x &  &  &  & x &  & outdoor (campus, urban) & ground (car) &  &  & x &  &  & x &  &  & x & x &  &  & x & x & x & x & RTK-GPS & LCM log & -- & -- & -- & 2m & --\\
\hline
\citetitle{dataset:utias} & 2011 & x &  &  &  &  &  & indoor (empty space) & ground (iRobot Create) & x &  & x & x &  &  &  &  &  &  &  &  &  &  & x &  & external tracking system & jpg (img), dat (non-img) & -- & -- & 4.78 & -- & 9\\
\hline
\citetitle{dataset:alderley} & 2012 & x &  & x &  & x &  & outdoor (urban) & ground (car) &  &  & x & x &  &  &  &  &  &  &  &  &  &  &  &  & manual & -- & 16 & 8 & -- & -- & 2\\
\hline
\citetitle{dataset:tum-rgbd} & 2012 & x &  &  &  & x &  & indoor (office, industrial hall) & ground (handheld, Pioneer 3) &  &  & x &  &  &  & x &  &  &  &  &  & x &  & x & x & external tracking system & plain text (non-img), png (img + depth), ROS bag & 0.285 & -- & 0.35 & -- & 15\\
\hline
\citetitle{dataset:cobot} & 2013 & x &  &  & x & x & x & indoor (office) & ground (robot) & x &  & x &  &  &  & x &  & x &  &  &  &  &  &  & x & -- & ROS bag & 131 & -- & 260 & 2y3m & 1082\\
\hline
\citetitle{dataset:kitti} & 2013 & x &  &  &  & x &  & outdoor (urban) & ground (car) &  & x & x & x & x &  &  &  &  & x &  &  & x & x & x & x & RTK-GPS & png (img), binary (laser), plain text (imu, gps) & -- & -- & 1.18 & 8d & 61\\
\hline
\citetitle{dataset:mit-stata} & 2013 & x &  &  & x & x & x & indoor (office) & ground (PR2) & x &  & x &  & x &  & x &  & x &  &  &  & x &  &  & x & map model & ROS bag & 42 & -- & 38 & 1y9m & 84\\
\hline
\citetitle{dataset:nordlandsbanen} & 2013 & x &  & x & x &  &  & outdoor (railway) & ground (train) &  &  & x & x &  &  &  &  &  &  &  &  &  & x &  &  & GPS & mp4 (video stream), plain text (gps) & 2916 & 729 & 39.74 & -- & 4\\
\hline
\citetitle{dataset:gardens-qut} & 2014 & x & x &  &  & x &  & indoor, outdoor (campus) & ground (handheld) &  &  & x & x &  &  &  &  &  &  &  &  &  &  &  &  & ground-plane position & png (imgs), plain text (ground plane) & -- & -- & -- & -- & 3\\
\hline
\citetitle{dataset:lcas-strands} & 2014 & x & x &  & x & x &  & indoor (office) & ground (SCITOS-G5) &  &  & x &  &  &  & x &  & x &  &  &  &  &  &  & x & -- & ROS bag & -- & -- & -- & 1y1m & 368\\
\hline
\citetitle{dataset:kaist} & 2015 & x & x &  &  & x &  & outdoor (urban) & ground (car) &  &  & x &  & x &  &  & x &  &  &  &  & x & x & x & x & RTK-GPS & png (imgs), plain text (imu, gps) & 84 & -- & -- & 18d & 36\\
\hline
\citetitle{dataset:euroc} & 2016 & x &  &  &  &  &  & indoor (industrial hall, office) & air (AscTec Firefly) &  & x &  & x & x &  &  &  &  &  &  &  & x &  & x & x & external tracking system & ROS bag & 0.8936 & -- & 0.37 & -- & 11\\
\hline
\citetitle{dataset:nclt} & 2016 & x &  & x & x & x &  & indoor, outdoor (campus) & ground (Segway) &  &  & x &  &  & x &  &  & x & x &  &  & x & x & x & x & RTK-GPS, SLAM-based & binary (laser), tiff (img), plain text (non-laser or img) & 147.4 & -- & 34.9 & 1y4m & 27\\
\hline
\citetitle{dataset:kudamm} & 2017 & x &  &  &  & x &  & outdoor (urban) & ground (car) &  &  & x & x &  &  &  &  &  &  &  &  &  &  &  &  & manual & jpg (img) & -- & -- & -- & -- & 2\\
\hline
\citetitle{dataset:oxford-robotcar} & 2017 & x & x & x & x & x &  & outdoor (urban) & ground (car) & x &  & x & x & x &  &  &  & x & x &  &  & x & x & x & x & RTK-GPS & png (img), binary (laser), plain text (imu, gps, odo) & 1010.46 & 10 & -- & 1y8m & 133\\
\hline
\citetitle{dataset:yq21} & 2017 & x &  &  &  & x &  & outdoor (campus) & ground (car) &  &  & x & x & x &  &  &  &  & x &  &  & x & x & x & x & RTK-GPS & binary (laser), jpg (imgs), plain text (gps) & 23 & -- & 6.5 & 1w & 21\\
\hline
\citetitle{dataset:cmu-seasons} & 2018 & x &  &  & x & x &  & outdoor (urban) & ground (car) &  &  & x & x &  &  &  &  &  &  &  &  &  &  & x & x & manual & jpg (img) & -- & 8.5 & -- & 330d & 17\\
\hline
\citetitle{dataset:fas} & 2018 & x &  &  & x & x &  & outdoor (urban) & ground (car) &  &  & x &  & x &  &  &  &  &  &  &  &  & x &  &  & GPS, manual & jpg (img) & 110 & -- & -- & 3y & 3\\
\hline
\citetitle{dataset:robotcar-seasons} & 2018 & x &  & x &  & x &  & outdoor (urban) & ground (car) &  &  & x & x & x &  &  &  &  &  &  &  &  &  & x & x & manual & jpg (img) & -- & 10 & -- & 178d & 10\\
\hline
\citetitle{dataset:bonn} & 2019 &  &  &  &  & x &  & indoor (office) & ground &  &  & x &  &  &  & x &  &  &  &  &  &  &  & x & x & external tracking system & png (imgs, depth), plain text (imu, gps) & -- & -- & -- & -- & 26\\
\hline
\citetitle{dataset:cbd} & 2019 & x &  &  &  & x &  & outdoor (urban) & ground &  &  & x & x & x &  &  &  &  &  &  &  &  &  &  &  & manual & png (imgs) & -- & -- & -- & -- & 1\\
\hline
\citetitle{dataset:mulran} & 2020 &  &  &  &  & x &  & outdoor (urban) & ground (car) &  &  &  &  &  &  &  &  &  & x & x &  & x & x &  & x & SLAM-based & binary (laser), CSV (global pose, radar ray), png (radar polar img) & 41.2 & -- & -- & 2m15d & 12\\
\hline
\citetitle{dataset:robotcar-radar} & 2020 & x &  & x &  & x &  & outdoor (urban) & ground (car) & x &  & x & x & x &  &  &  &  & x & x &  & x & x & x & x & RTK-GPS, SLAM-based & png (img, raw laser, radar), binary (laser), plain text (imu, gps, odo) & 280 & 10 & -- & 1m & 32\\
\hline
\citetitle{dataset:usyd} & 2020 & x &  & x & x & x &  & outdoor (campus) & ground (car) & x &  & x & x &  &  &  &  &  & x &  &  & x & x & x & x & GPS & ROS bag & -- & -- & -- & 1y & 52\\
\hline
\citetitle{dataset:iplt} & 2021 & x & x & x & x & x &  & outdoor (parking) & ground (car) & x & x &  & x &  &  &  &  & x &  &  &  &  & x & x & x & GPS & ROS bag & -- & 0.2 & -- & 2y & 127\\
\hline
\citetitle{dataset:radiate} & 2021 & x & x & x &  & x &  & outdoor (parking, urban) & ground (car) &  &  & x & x & x &  &  &  &  & x & x &  & x & x & x & x & RTK-GPS & ROS bag & -- & -- & 4.98 & -- & --\\
\hline
\citetitle{dataset:ntu} & 2022 & x &  &  &  &  &  & indoor, outdoor (campus) & air (DJI M600) &  & x &  & x & x &  &  &  &  & x &  &  & x &  & x & x & external tracking system & ROS bag & 1.845 & -- & 0.9 & -- & 9\\



%  \end{tabular}}
%\end{table*}
\end{longtable}

\end{tiny}

\twocolumn

\renewcommand{\arraystretch}{1.25}
\setlength{\tabcolsep}{3pt}


\subsection{Final observations}
\label{sec:discussion:observations}


\section{Challenges and Future Directions}
\label{sec:future}

\begin{itemize}[nosep]
\item multi-robot long-term localization and mapping
\item fusion of different sensors
\item availability of datasets in continuous operation in challenging long-term scenarios (urban, industry -- storage / logistics facilities)
\item edge/cloud computing
\item parameters of the algorithms adapting along time accordingly to changes in the environment
\item specific hardware for improving computational performance (FPGA, operating systems specific for real-time usage, etc.)
\item interaction with the user: map update and consequent visualization by the users not being on raw data but instead on higher levels of abstraction (CAD, automatic extract of the environment topological locations)
\item hierarchical localization and mapping (topological location > coarse localization > millimeter level of localization or  something like that)
\end{itemize}

\section{Limitations of the Study}
\label{sec:limitations}

Section to discuss possible limitations of the study (timeframe, wider approach, etc.).

\begin{itemize}[nosep]
\item only one query for discussion: e.g., for searching datasets, possibly, a different query should have been used
\item overview long-term SLAM vs in-depth analysis and discussion of each type of techniques: our review synthesizes all types of techniques, if the reader wants an in-depth analysis, different reviews should be performed
\item related to previous one, each category appearance dynamics sparsity multi-session and computational have all different aspects related to each other, that probably should be treated differently in different data extraction items to improve the data organization -- however the main goal of this review was to overview all trends and not focus on one specifically, so this categorization should be done separately and probably in singular literature reviews.
\item limited information on the experiments conditions (traveled distance, duration, etc.) given by the authors
\item some works such as PointNetVLAD, FAB-MAP and SeqSLAM (+LOAM, and LEGO-LOAM) did not appeared in the identification phase. However, the ones included have improvements over all of these...; also ``Robust optimization of factor graphs by using condensed measurements''
\item discussion does not focus on the datasets used but rather on the methodologies, to not extend even further the review
\item ORBSLAM2, ORBSLAM3 not identified
\item in retrospective, a bias of the authors may occurred relative and influenced in the inclusion of articles related to air-vehicles, underwater, and multi-robot systems -- air although batteries possibly a limitation their viewpoint specially on bird eye camera view is completely different than usual ground robots (facing forward) and so long-term methods for ground may not be as effective due to different viewpoint, underwater due to low number of articles found and the importance of long-term SLAM in these environments (plants growing, different conditions of the water in camera), and multi-robot systems due to the future these systems can improve greatly the mapping and planning efficiency of robotic installations
\end{itemize}

\section{Conclusions}
\label{sec:conclusions}

This paper presents a systematic literature review on long-term localization and mapping for mobile robots. The review selects 142 works from the literature covering the main strategies to achieve lifelong SLAM and discussing the experimental data (including private experiments and public datasets) and evaluation metrics commonly used to assess the performance of autonomous systems in long-term operations.
The discussion analyzes the included works in terms of appearance invariance to changing conditions in the environment, dealing with dynamic elements in the scene, multi-session strategies, map sparsification techniques to bound the computational resources, and other computational-related topics.
Also, an overview over the bibliographic data of the included works identifies the most used terms in long-term SLAM and identifying research networks between authors using the VOSviewer~\parencite{results:vosviewer:1,results:vosviewer:2} tool, while also discussing the evolution of the number of publications over time and the publication venues with more records.

Overall, the methodologies discussed in this study are a step forward to achieve lifelong autonomy.
In terms of dealing with appearance changes in the environment, CNN-based features are more discriminative compared to handcrafted features, while considering both appearance and geometric cues in the descriptor improve its invariance to changing conditions. Semantic features and considering different sensorization sources (e.g., cameras, LiDAR, and/or radar) can also increase the robustness to appearance changes in the environment.
For dealing with dynamic elements, the most common approach is to distinguish between static, dynamic, and semi-static changes in the perceived environment, while only using static permanent changes for localization to improve the reliability of the pose estimator.
As for constraining the map size to the explored environment area instead of the trajectory length, the evaluation of the mutual information by using techniques based on information theory or heuristics is the most currently used technique to bound the map size.

However, there are still challenges and future directions for the research on long-term SLAM. Vision-based global recognition is affected by viewpoint variance, where using omnidirectional sensors to increase the field-of-view relative to perspective cameras may decrease that variance. Decentralized computation architectures may be an interesting research direction to offload heavier computation tasks from the mapping agent and improve the scalability of long-term SLAM algorithms. For example, cloud and edge computing devices allow the parallelization of tasks and increase the availability of computational resources, even though communication constraints should be accounted. Also, the extension of the methodologies discussed in this study to multi-robots systems may optimize the operation of autonomous systems while also increasing their robustness in the long-term.

Lastly, this review can be updated in future studies thanks to following a systematic review that defines explicitly the time interval in the selection process (the data of the last full inquiry is May 17, 2022) and the search query used to identify records from the literature.
Also, all documentation and scripts using during the review process are available in a public GitHub repository to facilitate replicating of the results and future updates to this review.
This paper may be extended to focus on single challenges of long-term SLAM for providing an in-depth comparison with experimental results between different methodologies.


\section*{Acknowledgments}

This study was supported by FCT -- Fundação para a Ciência e a Tecnologia and INESC TEC -- Institute for Systems and Computer Engineering, Technology and Science.

\section*{ORCID}

\noindent \textit{Ricardo B. Sousa} \includegraphics[width=1em]{orcid.pdf} \href{https://orcid.org/0000-0003-4537-5095}{https://orcid.org/0000-0003-4537-5095}

\noindent \textit{H\'{e}ber M. Sobreira} \includegraphics[width=1em]{orcid.pdf} \href{https://orcid.org/0000-0002-8055-1093}{https://orcid.org/0000-0002-8055-1093}

\noindent \textit{Ant\'{o}nio Paulo Moreira} \includegraphics[width=1em]{orcid.pdf} \href{https://orcid.org/0000-0001-8573-3147}{https://orcid.org/0000-0001-8573-3147}



%% BIBLIOGRAPHIC REFERENCES

\printbibliography

%\bibliographystyle{unsrt}
%\bibliography{references.bib}



%% AUTHORS BIOGRAPHIES

\newpage

%\hfill\break

\begin{tabular}{p{0.25\columnwidth} p{0.68\columnwidth}}
\raisebox{-0.925\totalheight}{\includegraphics[width=0.25\columnwidth]{figures/sousa.jpg}} &
\textbf{Ricardo B. Sousa} obtained a Master of Science (M.Sc.) degree in Electric and Computers Engineering (ECE) at Faculty of Engineering of the University of Porto (FEUP), in 2020.
He is currently working towards the Ph.D. degree in electrical and computer engineering with FEUP, and he has a graduate research scholarship from FCT -- Fundação para a Ciência e a Tecnologia at the Centre for Robotics in Industry and Intelligent Systems from INESC TEC.
Also, he is an invited assistant lecturing the courses Software Design and Industrial Informatics from the M.Sc. in ECE at FEUP.
His research interests include robotics, sensor fusion, and localization and mapping for autonomous robots.\\
\end{tabular}

%\hfill\break
%\newpage

\begin{tabular}{p{0.25\columnwidth} p{0.68\columnwidth}}
\raisebox{-0.925\totalheight}{\includegraphics[width=0.25\columnwidth]{figures/sobreira.jpg}} &
\textbf{H\'{e}ber M. Sobreira} was born in Leiria, Portugal, in July 1985. He graduated with an M.Sc. degree (2009) and a Ph.D. degree (2017) in Electrical Engineering from the University of Porto. Since 2009, he has been developing his research within the Centre for Robotics in Industry and Intelligent Systems at INESC TEC. His main research areas are navigation and control of indoor autonomous vehicles.\\
\end{tabular}

%\hfill\break
%\newpage

\begin{tabular}{p{0.25\columnwidth} p{0.68\columnwidth}}
\raisebox{-0.925\totalheight}{\includegraphics[width=0.25\columnwidth]{figures/moreira.png}} &
\textbf{Ant\'{o}nio Paulo Moreira} graduated with a degree in electrical engineering at the University of Oporto, in 1986. Then, he pursued graduate studies at University of Porto, obtaining a M.Sc. degree in electrical engineering -- systems in 1991 and a Ph.D. degree in electrical engineering in 1998. Presently, he is Associate Professor with tenure at the Faculty of Engineering of the University of Porto and researcher and head of the Centre for Robotics in Industry and Intelligent Systems at INESC TEC. His main research interests are process control and robotics.\\
\end{tabular}



%% APPENDIXES

\cleardoublepage

\section{Introduction}
\label{sec:intro}

An autonomous mobile robot requires a representation of its surroundings to localize itself relative to the environment.
Simultaneous Localization and Mapping (SLAM) addresses this problem by incorporating the robot state estimation (pose and possibly other state variables) concurrently with the mapping process. This process builds a representation of the environment perceived by the robot originating a map incrementally built when exploring unknown areas or refined on passages through known locations.

In a static scene, the robot would only need to map once because it would be always consistent with the environment.
However, autonomous systems deployed in industrial locations, outdoor environments, or even service-oriented applications such as shopping centers or homes deal with moving elements in the scene (humans, objects), environment reconfiguration (machines moving places), and appearance variations (lighting, weather, seasonal, or daytime changes).
These varying conditions are a challenge for the SLAM system, where the system should decide how (consider the most current state or only the most static changes?) and when to update the map (when the variations occur or after a certain time?).
This challenge is also known as the stability-plasticity dilemma, where the long-term localization and mapping should both adapt to new environment changes and preserve old states over time~\parencite{biber-duckett:2009:0278364908096286}.

Furthermore, the map would grow indefinitely when gathering new information from the environment. This ever-growing problem poses another challenge for the SLAM system in the long-term due to the limited computational resources of the mapping vehicle. Indeed, the map size should be dependent on the environment area and not on the trajectory length of the mapping process nor operation time, only growing when the robot would explore unknown locations~\parencite{kretzschmar-stachniss:2012:0278364912455072}.

Although several studies overview SLAM literature, only a subset of those studies mentions long-term challenges of performing SLAM.
\cite{purpose:study:cadena:2016} has a brief survey on the robustness and scalability of autonomous systems focused on loop closure validation, dynamic environments, pose graph sparsification, and parallel and distributed computing for metric and semantic SLAM.
In contrast, \cite{purpose:study:lowry:2016} limits the study to vision-based topological SLAM discussing also the challenge of varying conditions.
While \cite{purpose:study:bresson:2017} overviews autonomous driving trends in terms of scalability, map updatability, and dynamicity, the survey limits the discussion to algorithms that have both odometry and mapping modules, excluding localization-only works. Also, \cite{purpose:study:bresson:2017} focuses more on the modules of the SLAM (relocalization, localization against a map).
\cite{purpose:study:kunze:2018} gives a brief overview of AI-related works for robustness to appearance changes and learning dynamics of the environment, discussing areas in which artificial intelligence (AI) enables long-term operation of autonomous systems.
As for \cite{purpose:study:zaffar:2018}, the study evaluates the long-term autonomy of sensors such as monocular and stereo cameras, and LiDAR.
However, to the best of the authors knowledge, none of the existing studies overviews the trends for dealing with long-term challenges in SLAM. Also, the studies that overview some of the challenges of lifelong SLAM do not clarify the process for identifying the cited works, not allowing other researchers to repeat the identification process of records when searching in bibliographic databases.

Therefore, this paper presents a systematic literature review on long-term localization and mapping following the Preferred Reporting Items for Systematic reviews and Meta-Analysis~(PRISMA)~\parencite{methodology:prisma} statement.
The systematic method followed in this review allows the replication of the results by other researchers and leads to the inclusion of 142 works for discussion and analysis.
Also, this paper makes available a public GitHub repository\footnote{\url{https://github.com/sousarbarb/slr-ltlm-mr}} with all the documentation and scripts used during the process of systematic revision of the literature.
The main contributions of this paper relative to the existing studies on SLAM are the following ones:

\begin{itemize}[nosep]
\item discussion on methodologies and trends focused on appearance invariance, dynamic elements, map sparsification and multi-session techniques, and other computational concerns;
\item comparative analysis on the public datasets and experimental data used in the included works in terms of environment conditions, sensorization, and distance and time properties;
\item presentation of common evaluation metrics used by the included works in the experimental results.
\end{itemize}

This review does not intend to review the fundamentals of SLAM nor the main formulations. The reader should refer to \cite{purpose:study:durrant-whyte-bailey:2006:1} and \cite{purpose:study:durrant-whyte-bailey:2006:2} for the Extended Kalman Filter (EKF) and particle filter probabilistic formulations, and to \cite{purpose:study:grisetti:2010} for the pose graph formulation of SLAM.





\subsection{Paper organization}

The rest of this review is organized as follows.
Section~\ref{sec:purpose} discusses the limitations of existent studies and reviews and presents the purpose and motivations of this paper.
Section~\ref{sec:methodology} explains the methodology followed in this review to search and select the included records, and the data extraction process for synthesis and analysis. In Appendix~\ref{a2:data-extraction}, Table~\ref{tab:a2:data-extraction} presents the data extraction results of the included records.
Next, Section~\ref{sec:overview} analyzes the bibliographic information of the 142 included works in this review in terms of the identification results of each data source considered in the methodology, keywords co-occurrence and co-authorship relations, the year of publication, and the publication venue.
Section~\ref{sec:discussion} discusses the methodologies found in the included records related to long-term localization and mapping and analyzes the experimental data and evaluation metrics used in the experiments by the authors.
Then, Section~\ref{sec:future} outlines challenges and future directions.
Section~\ref{sec:limitations} discusses possible limitations of this study.
Lastly, Section~\ref{sec:conclusions} presents the conclusions.

\section{Purpose of the study}
\label{sec:purpose}

\textcolor{red}{PUT INTRO HERE!!!}





\subsection{Limitations of current studies}



\begin{table}[h]
  \caption[Existent Literature Reviews, Surveys, and Tutorials on SLAM.]{Existent Literature Reviews, Surveys, and Tutorials on SLAM.}
  \label{tab:purpose:current-literature}
  \centering
  {\scriptsize
  \begin{tabular}{c p{0.43\columnwidth} p{0.43\columnwidth}}

\hline
\textbf{Year} & \textbf{Topic} & \textbf{Reference}\\
\hline
2006 & Probabilistic formulations (EKF, particle filter)
     & \cite{purpose:study:durrant-whyte-bailey:2006:1}, \cite{purpose:study:durrant-whyte-bailey:2006:2}\\
\hline
2008 & Probabilistic and pose graph formulations
     & \cite{purpose:study:thrun:2008}\\
\hline
2010 & GraphSLAM
     & \cite{purpose:study:grisetti:2010}\\
\hline
2011 & Observability, convergence, consistency (feature-based SLAM)
     & \cite{purpose:study:dissanayake:2011}\\
\hline
2012 & Visual odometry
     & \cite{purpose:study:scaramuzza-fraundorfer:2011:1}, \cite{purpose:study:scaramuzza-fraundorfer:2012:2}\\
\hline
2014 & Underwater navigation and localization
     & \cite{purpose:study:paull:2014}\\
\hline
2015 & Visual SLAM
     & \cite{purpose:study:yousif:2015}\\
\hline
2016 & Observability, convergence, consistency (feature and graph-based SLAM)
     & \cite{purpose:study:huang-dissanayake:2016}\\
     & Multi-robot SLAM
     & \cite{purpose:study:saeedi:2016}\\
     & Visual place recognition
     & \cite{purpose:study:lowry:2016}\\
     & SLAM literature overview
     & \cite{purpose:study:cadena:2016}\\
\hline
2017 & Autonomous vehicles
     & \cite{purpose:study:bresson:2017}\\
\cline{1-1}
\cline{3-3}
2018 &
     & \cite{purpose:study:kuutti:2018}\\
\hline
2018 & AI for long-term autonomy
     & \cite{purpose:study:kunze:2018}\\
     & Long-term sensorization
     & \cite{purpose:study:zaffar:2018}\\
\hline
2020 & Deep learning
     & \cite{purpose:study:fayyad:2020}\\
     & Multi-robot search and rescue
     & \cite{purpose:study:queralta:2020}\\
\hline
2021 & Self-driving vehicles
     & \cite{purpose:study:badue:2021}\\
\hline
2022 & Underground navigation
     & \cite{purpose:study:ebadi}\\
\hline

  \end{tabular}}
\end{table}





\subsection{Motivations and goals}

Research question: What is the current state of the art of long-term localization and mapping using mobile robots?

Goals of this review:

\begin{itemize}[nosep]
\item which are the main strategies for accomplishing long-term operations with mobile robots;
\item how to deal with varying conditions of the environment;
\item how do autonomous robots deal with the dynamics of the environment;
\item which are the main strategies to deal with the limited computational resources of a mobile robot on long-term operations;
\item how the methods evaluate their results;
\item which are the public datasets more used for evaluating long-term localization and mapping.
\end{itemize}

PICO framework (Population--Intervention--Comparison--Outcome) helps to frame the research questions of this systematic review into a more structured framework:

\begin{itemize}[nosep]
\item \textbf{Population:} mobile robots;
\item \textbf{Intervention:} localization, mapping, SLAM;
\item \textbf{Comparison:} \textit{not applicable to this study};
\item \textbf{Outcome:} long-term operation, lifelong autonomy, robust.
%\item \textbf{Context:} continuous operation, service robots, industrial environments.
\end{itemize}


\section{Methodology}
\label{sec:methodology}

A systematic literature review uses explicit, rigorous, and reproducible systematic methods to synthesize the findings of studies related to a particular research question, topic area, or phenomenon of interest. This type of review assures the quality and trustworthiness of the review's findings by presenting a complete, organized, and summarized analysis of all works considered while allowing others to replicate or update the reviews. The most common standard for performing a systematic review is the PRISMA~\parencite{methodology:prisma} statement. Although the PRISMA statement has been designed originally for evaluating the effects of health interventions, the checklist items of the methodology are general and applicable to other subject areas. Thus, the methodology used in this systematic review follows the PRISMA~\parencite{methodology:prisma} guidelines.

This section presents the detailed methodology used in this study. First, the eligibility criteria decide which studies to include in the review. Next, the search strategy details the information sources considered in the review and the base string and search fields used for inquiring these sources. Furthermore, the selection process focuses on describing its stages and the quality evaluation criteria used to select works for the synthesis and analysis phase of the review. Lastly, the data extraction process details the relevant data collected for synthesis and analysis.
Parsifal \parencite{methodology:parsifal} is the online tool used to support the literature review in designing the methodology protocol, removing duplicates, screening and selecting works including their quality assessment. Additional documentation and scripts developed within the scope of this review related to removing duplicates, checking and processing the bibliographic references, and data extraction are available in the public GitHub repository.

\subsection{Eligibility criteria}
\label{sec:methodology:eligible}

Table~\ref{tab:methodology:exclusion-criteria} presents the exclusion criteria used to determine the eligible studies for the selection process. These eligibility criteria focus mainly on the type of paper and availability. The index criterion rejects all publications not indexed in a scientific publication venue. This rejection guarantees that the eligible works were peer-reviewed by the scientific community. Also, the exclusion criteria reject short papers and gray, secondary, and tertiary literature. Short papers do not usually present a detailed methodology of their scientific contribution. As for only considering primary literature in the review, this criterion increases the relevance of search results by favoring original articles and simultaneously guaranteeing peer-revision of the works. In terms of language, only considering studies with English full-texts increases the scope and visibility of the review. Similarly, the eligibility criteria reject studies not available in digital libraries for reproducibility and accessibility reasons.

\begin{table}[h]
  \renewcommand{\arraystretch}{1.25}
  \setlength{\tabcolsep}{3pt}
  \caption[Exclusion criteria for the selection process.]{Exclusion criteria for the selection process.}
  \label{tab:methodology:exclusion-criteria}
  \centering
  {\scriptsize
  \begin{tabular}{c m{0.28\columnwidth} m{0.6\columnwidth}}

\hline
\textbf{E\#} & \textbf{Criteria} & \textbf{Statement}\\
\hline
E1 &
Index &
Papers not indexed in a scientific publication venue\\
\hline
E2 &
Language &
Full-text of the papers not published in English\\
\hline
E3 &
Subject Area &
Papers not classified in the databases as Computer Science, Engineering, Mathematics, or Multidisciplinary\\
\hline
E4 &
Short Papers &
Papers classified as short papers accordingly to the publication venue\\
\hline
E5 &
Gray, Secondary, and Tertiary Literature &
Books, preprints, reports, reviews, thesis, ...\\
\hline
E6 &
Availability &
Full-text of the papers not available in digital libraries\\
\hline
E7 &
Dataset &
Papers that focus only on data collection\\
\hline
E8 &
Coverage &
Papers using only odometry for localization\\
\hline
E9 &
Scope &
Papers that focus on different and not related subjects\\
\hline

  \end{tabular}}
\end{table}

Another exclusion criterion considered in the review is relative to the studies' categorization of their subject areas by bibliographic databases. The ones considered in the review are Computer Science, Engineering, Mathematics, or Multidisciplinary areas. In the list provided by the Clarivate's Journal Citation Reports\footnote{\url{https://jcr.clarivate.com/jcr/browse-categories}}, these four subject areas include the artificial intelligence, interdisciplinary applications, electrical and computers engineering, robotics, and applied mathematics categories, among others. These categories are intrinsically related to the localization and mapping problem for long-term operation of mobile robots.

The final three criteria presented in Table~\ref{tab:methodology:exclusion-criteria} focus on the scientific contribution of the studies. The dataset criterion rejects all works that focus only on sharing a data collection. Although these works are important for the evolution of localization and mapping algorithms in providing a benchmark for comparison and reference purposes, their scientific contribution is not directly comparable to research articles. Odometry-only approaches are unusable over long distances invalidating their use for long-term operations with mobile robots. As for the scope criterion, this review does not consider as eligible papers not related to long-term localization and mapping that do not fall under the other exclusion criteria.

\subsection{Search strategy}
\label{sec:methodology:search}

The search phase consists of identifying the data sources that could be relevant for this literature review, and defining the base string and which search fields considered to obtain the results for the review. \citetitle{methodology:search:db:wos} and \citetitle{methodology:search:db:scopus} are traditionally the two most widely used bibliographic databases. However, previous studies demonstrate that different databases differ significantly in their scientific coverage~\parencite{methodology:search:db:coverage:scopus-wos,methodology:search:db:coverage:dim-scopus-wos}. Thus, the data sources considered in this review are the following ones: \citetitle{methodology:search:db:acm}, \citetitle{methodology:search:db:dimensions}, \citetitle{methodology:search:db:ieee-xplore}, \citetitle{methodology:search:db:inspec}, \citetitle{methodology:search:db:scopus}, and \citetitle{methodology:search:db:wos}.

Moreover, May 17, 2022, is the date of the last full inquiry. Future reviews on the topic of this study should consider this final date as theirs initial one. As for inquiring the data sources, the base string used is the following one:

\vspace{1em}

\noindent\begin{center}
\begin{tabular}{l l}
\texttt{(robot* OR vehicle*)}                     & \texttt{AND}\\
\texttt{((locali* AND map*) OR "slam")}           & \texttt{AND}\\
\texttt{("long term" OR "life long" OR lifelong)}\\
\end{tabular}
\end{center}

\vspace{1em}

The first terms, \texttt{robot* OR vehicle*}, attempt to focus the search results to the desired population. These two terms have multiple synonyms within the scope of autonomous mobile robots: mobile robots, autonomous vehicles, robotics, agricultural robots, intelligent robots, service robots, unmanned aerial/ground/underwater vehicles, among other terms. Therefore, by adding the asterisk to the end of the terms robot and vehicle (\texttt{robot*} and \texttt{vehicle*}, respectively), and by only considering the terms with asterisk in the inquiry, all the synonyms are covered for the desired population. Given the incompatibility of the \citetitle{methodology:search:db:dimensions} database with wildcards (e.g., using the asterisk), the first part of the base string becomes as follows when searching in this database: \texttt{robot OR robots OR robotics OR vehicle OR vehicles}.

The next part of the query focus on the intervention side of the systematic review. Given the interest of this review on searching for localization and mapping algorithms, \texttt{locali*} and \texttt{map*} summarize all the synonyms for the localization and mapping terms, respectively. For example, \texttt{locali*} not only is agnostic to the US versus UK spelling differences (localization vs localisation, respectively) but also resumes several synonyms: localization, localize, or localizing. The term \texttt{map*} also attempts to cover its respective synonyms such as map, maps, or mapping.
Also, the acronym \texttt{"slam"} is another alternative to search for localization and mapping algorithms. Even though its definition is compatible with \texttt{locali* AND map*}, some authors only refer to SLAM.
Similarly to the inquiry's first part, the second one becomes as follows for searching in \citetitle{methodology:search:db:dimensions}: \texttt{((localize OR localization OR localizing OR localise OR localisation OR localising) AND (map OR maps OR mapping)) OR "slam"}.

As for \texttt{"long term" OR "life long" OR lifelong}, this part of the base string is relative to the outcome of the PICO framework, presented in Section~\ref{sec:purpose}. The reason for having both \texttt{"life long"} and \texttt{lifelong} terms is the existing confusion in which term is grammatically the correct one.

Furthermore, the Title, Abstract, and Keywords are the fields considered for obtaining the search results. The third one includes the author keywords, the indexed terms by the databases, and the uncontrolled ones if they are available. The selection of these search fields for this review improves the relevance of the results compared to using all fields and the full text by focusing the search on the summary items of the works. Indeed, the main contributions of scientific works should be summarized in at least the title, abstract, or the author keywords. The indexed terms also help in obtaining records only related to the base string used in this review.
However, not all data sources have available the search fields considered in the review or some of them require an adaptation when performing the search. 
Although the \citetitle{methodology:search:db:acm} allows searching within multiple search fields, including the ones considered in this review, the advanced search query on this library sets by default an AND operator between the different fields. This setting must be changed manually in the query syntax to the desired OR operator. Also, there are two options to search items in the \citetitle{methodology:search:db:acm}: \textit{The ACM Full-Text Collection} and \textit{The ACM Guide to Computing Literature}. Given that the latter includes all the content from the former, the identification process in this source performs the search using \textit{The ACM Guide to Computing Literature} option.
Other than searching in the publications' full data, \citetitle{methodology:search:db:dimensions} only has the title and abstract search fields compatible with this review. Given the limitation of \citetitle{methodology:search:db:ieee-xplore} to 7 wildcards, the search results of this digital library using the base string for the inquiry are the grouping of different searches considering only a search field at a time, importing each search results to Parsifal and removing the duplicates. As for \citetitle{methodology:search:db:inspec}, \citetitle{methodology:search:db:scopus}, and \citetitle{methodology:search:db:wos}, these databases have available all the search fields considered in the review.

In terms of the publication date, this review does not restrict it to avoid ignoring important works and to improve the discussion. Indeed, to best of the authors knowledge, there is not available a systematic review on long-term localization and mapping for mobile robots to provide an initial date for rejecting older publications. Even though the number of publications per year could indicate an initial date on when the topic gained relevance, the date filtering could still reject important works.

\subsection{Selection process}
\label{sec:methodology:selection}

The selection process of this review summarized in Figure~\ref{fig:methodology:prisma-flow} has three phases: identification, screening, and quality assessment. The first phase consists of inquiring each data source discussed previously with the base string and adapting it if needed. The second phase requires screening the papers. In this review, screening is equivalent to reading the publications' title and abstract and deciding whether the study is eligible or not based on the exclusion criteria. Then, a set of evaluation criteria assesses the quality of the eligible records. The records obtained after the three phases of the selection process are for the data extraction phase.

\begin{figure}[h]
  \centering
  \includegraphics[width=\columnwidth]{figures/selection.pdf}
  \caption{PRISMA flow diagram for the selection process.}
  \label{fig:methodology:prisma-flow}
\end{figure}

\subsubsection{Identification}

In the identification phase of this review, the search strategy is applied to all data sources. \citetitle{methodology:search:db:acm}, \citetitle{methodology:search:db:dimensions}, \citetitle{methodology:search:db:inspec}, \citetitle{methodology:search:db:scopus}, and \citetitle{methodology:search:db:wos} data sources only require a single inquiry to obtain the search results.
Given the limitation of the \citetitle{methodology:search:db:ieee-xplore} for using wildcards mentioned in Section~\ref{sec:methodology:search}, the number of records for this source presented in Figure~\ref{fig:methodology:prisma-flow} represents the results of 7 inquiries (using the fields title, abstract, author keywords, IEEE terms, INSPEC controlled terms, and the INSPEC uncontrolled ones, respectively) after removing the duplicates with the support of Parsifal. Although the total number of search results found is 2160, Parsifal is used to remove duplicates from different data sources, excluding 1339 records. Following the duplicates removal, the exclusion criteria defined in Section~\ref{sec:methodology:search} exclude 232 works from the review. This exclusion is possible due to \citetitle{methodology:search:db:inspec}, \citetitle{methodology:search:db:scopus}, or \citetitle{methodology:search:db:wos} having filters related to the publication's type, subject area, and language.

The works excluded from the search results also include the ones that do not meet the exclusion criteria E4 and E7. For the first one, a Python script available in the GitHub repository of this review searches studies with a number of pages lower or equal to 4. Even though short papers have a maximum number of 3 pages, the papers with 4 pages do not usually present a detailed methodology.
As for the E7 exclusion criterion, some works are possible to remove from the review by searching in their title for the term ``dataset''.
All excluded articles of this review are double-checked to certify if the exclusion criteria are correctly applied. For example, articles published in the Remote Sensing journal from MDPI do not meet the E3 criterion. Indeed, the Journal Citations Reports from Clarivate classifies it by the following categories: Remote Sensing, Geosciences Multidisciplinary, Environmental Sciences, and Imaging Science \& Photographic Technology. However, most search results from this journal found in the identification phase are directly related to the topic of this review and the respective subject areas. Thus, in these cases and in other ones related to the remaining exclusion criteria, the decision is reverted to consider the initially rejected studies for the next phase of the review.

\subsubsection{Screening}

Next, the screening phase in this review consists of reading the title and abstract of the publications and rejecting the ones that meet the exclusion criteria. However, the initially rejected papers have another assessment for validating the exclusion. The analysis of the results and conclusions of these publications considering the exclusion criteria either confirms the exclusion decision or reverses it to eligible works for quality assessment. As a result of the screening phase, 178 studies are rejected from the initial identified 589 works. The duplicate records found in screening and removed manually are due to titles with invalid characters originated by exporting the search results from the \citetitle{methodology:search:db:dimensions}~database.

\subsubsection{Quality assessment}

The quality evaluation in this review of the selected works from screening follows the 8 Quality Evaluation (QE) criteria presented in Table~\ref{tab:methodology:quality-assessment}. All of them are subjective criteria derived from the analysis of the eligible works. The score column establishes the possible values for the QE criteria, in which the minimum, intermediate, and maximum values correspond to none, partial, and full compliance, respectively.
Furthermore, QE1, QE2, QE4, and QE8 focus on the details provided in the papers, specifically, if the discussion of the related work, the proposed methodology, the experimental setup, and the results are detailed and thoroughly analyzed in the publication, respectively.
The possible scores for QE3 are twice the value of QE1, QE2, QE4, and QE8 due to this criterion being directly related to the topic of the review. A work focusing on both localization and mapping problems will have a score of 2.0 (full compliance). If the study only focuses on one of these problems or none of them, the scores will be 1.0 or 0.0, i.e., partial or no compliance, respectively.
QE5 evaluates the long-term results of the eligible studies and is either 2.0 (full) or 0.0 (no compliance). This criterion has the same range as QE3 for similar reasons, given the focus of this review on long-term localization and mapping algorithms.
The definition of long-term experiments for assigning full compliance in QE5 is the following one: dynamic changing environments (e.g., dynamic elements or semi-static ones), increasing environments or feature maps in terms of their size, redundant data removal, or varying conditions (e.g., different seasons of the year or lighting conditions).
QE6 and QE7 can only be 1.0 or 0.0. The former criterion intends to highlight works that compare themselves to the state of the art and/or ground-truth data. The latter emphasizes the importance of having available either the implementation of the proposed methodology or the data used in the experiments for other works to be able to compare the proposed methodologies.
Lastly, considering the possible scores for the QE criteria in Table~\ref{tab:methodology:quality-assessment}, each work can only have a  maximum score of 10.0.

\begin{table}[h]
  \renewcommand{\arraystretch}{1.25}
  \setlength{\tabcolsep}{3pt}
  \caption[Quality evaluation criteria and score range.]{Quality evaluation criteria and score range.}
  \label{tab:methodology:quality-assessment}
  \centering
  {\scriptsize
  \begin{tabular}{c m{0.68\columnwidth} c}

\hline
\textbf{QE\#} & \textbf{Criteria} & \textbf{Score}\\
\hline
QE1 &
Does the paper have an updated state of the art on long-term localization and mapping? &
\{0.0, 0.5, 1.0\}\\
\hline
QE2 &
Is the methodology appropriate and detailed? &
\{0.0, 0.5, 1.0\}\\
\hline
QE3 &
Does the methodology consider both localization and mapping problems? &
\{0.0, 1.0, 2.0\}\\
\hline
QE4 &
Is the hardware and/or software used in the experiments detailed? &
\{0.0, 0.5, 1.0\}\\
\hline
QE5 &
Does the paper presents any kind of long-term experimental results? &
\{0.0, 2.0\}\\
\hline
QE6 &
Does the paper presents comparative results with other methods and/or ground-truth data? &
\{0.0, 1.0\}\\
\hline
QE7 &
Does the work's implementation and/or the data used in the experiments are publicly available? &
\{0.0, 1.0\}\\
\hline
QE8 &
Is the discussion of the results and conclusions appropriate and detailed? &
\{0.0, 0.5, 1.0\}\\
\hline

  \end{tabular}}
\end{table}

After evaluating the 411 eligible works accordingly to the previously discussed QE criteria (the scores of each record are available in the GitHub repository), the first conclusion of the authors is that works with a non-detailed or not appropriate methodology, results' discussion, or conclusions should not be included in the review. Another conclusion is relative to rejecting works that do not consider either localization or mapping problems, or do not present any long-term experimental results, given the focus of this review on the long-term localization and mapping problem for mobile robots. Furthermore, the quality assessment phase should consider a cut-off score to filter works with low quality scores. Consequently, the assessment phase considers the following two reasons to reject a record:

\begin{enumerate}[nosep]
\item QE2, QE3, QE5, QE8: reject works with a 0.0 (no compliance) score;
\item cut-off score: reject works with a score lower or equal to 7.5/10.0.
\end{enumerate}

The distribution of the evaluation scores and the QE criteria itself justify the selection of a 7.5/10.0 cut-off score. 
Figure~\ref{fig:methodology:qe} illustrates the scores distribution for all eligible works versus the scores of the ones that pass the first criterion defined previously for the QE phase (related to the compliance on the QE2, QE3, QE5, and QE8 criteria). The assessment of this criterion rejects 116 records (28\%) of the 411 eligible works (see Figure~\ref{fig:methodology:prisma-flow}). 
Even though the distribution of the evaluation scores changes significantly in the range of scores lower or equal to 7.5/10.0, as observed between Figures~\ref{fig:methodology:qe:qe_wo-r1} and~\ref{fig:methodology:qe:qe}, only one work with a score higher than 7.5 is rejected due to not having a detailed and appropriate discussion of the results. This result indicates that interesting works are associated with high scores, as intended when using a quality assessment methodology, while also suggests that the range between 8.0 and 10.0 have the most interesting and quality works compatible with the focus of this review on long-term localization and mapping. Although only assessing the eligible works would seem to lead to the same results in terms of records included in the review, the rejection criterion on QE2/3/5/8 prevents outliers related to the quality assessment.
From the remaining 295 eligible works, cut-off scores from 7.5 up to 8.5 have the following corresponding rejection rates:

\begin{itemize}[nosep]
\item 7.5/10.0 \phantom{$\xrightarrow{\text{reject}}$} 120 records (40.7\%) \phantom{$\xrightarrow{\text{include}}$} 175 records
\item 8.0/10.0 $\xrightarrow{\text{reject}}$ 160 records (54.2\%) $\xrightarrow{\text{include}}$ 135 records
\item 8.5/10.0 \phantom{$\xrightarrow{\text{reject}}$} 203 records (68.8\%) \phantom{$\xrightarrow{\text{include}}$} 92 records
\end{itemize}

\begin{figure}[h]
  \centering
  \subfloat[][]{\includegraphics[width=\columnwidth]{figures/qe.png}%
  \label{fig:methodology:qe:qe}}
  \linebreak
  \subfloat[][]{\includegraphics[width=\columnwidth]{figures/qe_wo-r1.png}%
  \label{fig:methodology:qe:qe_wo-r1}}
  \caption{Distribution of the quality evaluation scores obtained from assessing the eligible works considered in the review: (a) all eligible works; (b) works that pass the rejection criterion during the QE assessment related to QE2/3/5/8 = 0.0 (no compliance).}
  \label{fig:methodology:qe}
\end{figure}

The 8.5 cut-off score would not be suitable because methods that focus only on localization or mapping, or not having either the implementation or the experimental data publicly available would be obligated to have maximum scores in the other criteria to be included in the review. In these cases, a work would have a maximum score of 9.0 due to partial compliance on QE3 or no compliance on the QE7 criteria. Likewise, a cut-off score of 8.0 would only leave a margin for having a single partial compliance on QE1, QE2, QE4 or QE8 criteria in similar cases, even though it would reject 160/295 (54\%) records. Therefore, the 7.5/10.0 cut-off score is more appropriate for the quality assessment phase in this review by leaving margin for works to have partial compliance in more than one criterion. Indeed, this cut-off score allows an article with no public data and/or implementation (e.g., due to confidentiality agreements) to have up to four criteria with partial compliance, depending on the criterion's maximum score or if the work has available the experiments data and/or implementation. Another example is articles that only focus on localization or mapping. In these cases, the work could have no public implementation, even though requiring a maximum score on all other criteria, or, if the work has public data or implementation available, two other criteria could have partial compliance.

Overall, as illustrated in Figure~\ref{fig:methodology:prisma-flow}, the quality assessment of the 411 eligible works considering the two rejection criteria previously mentioned leads to rejecting a total of 236 (57\%) records. As a result, the remaining 175 records will be analyzed for data extraction.

\subsection{Data extraction}
\label{sec:methodology:data}

The data extraction process analyzes the records selected after the quality assessment phase and extracts information from these works. In the scope of this review, the Data Extraction (DE) items required for each record are the following ones:

\begin{itemize}[nosep]
\item $\left[\textbf{DE1}\right]$ \textbf{Long-term considerations} -- long-term factors the works consider in their proposed approach and experiments. Considering the knowledge obtained in the previous phases of this review's methodology, the authors considered the following factors for categorizing the included works:
  \begin{itemize}[nosep]
  \item appearance: varying conditions, appearance changes;
  \item dynamics: environment dynamics, dynamic elements;
  \item sparsity: map pruning, redundant data removal;
  \item multi-session: map management;
  \item computational: memory management, efficiency.
  \end{itemize}
\item $\left[\textbf{DE2}\right]$ \textbf{Localization} -- how the robot localizes itself and the type of localizer;
\item $\left[\textbf{DE3}\right]$ \textbf{Mapping} -- type of the map;
\item $\left[\textbf{DE4}\right]$ \textbf{Multi-robot} -- if the proposed methodologies consider multi-robot systems;
\item $\left[\textbf{DE5}\right]$ \textbf{Execution mode} -- offline, online, if requires both, or if no information on this item;
\item $\left[\textbf{DE6}\right]$ \textbf{Environment and domain} -- type of environment (indoor, outdoor) and domains (air, ground, water) tested with the proposed methodologies;
\item $\left[\textbf{DE7}\right]$ \textbf{Sensory setup} -- which sensors considered in the methodologies;
\item $\left[\textbf{DE8}\right]$ \textbf{Non-public experiments} -- if the authors performed experiments or tests with non-public data;
\item $\left[\textbf{DE9}\right]$ \textbf{Ground-truth} -- how ground-truth for non-public data is obtained or its type, if available;
\item $\left[\textbf{DE10}\right]$ \textbf{Distance and time characteristics} -- relative to the non-public experiments if available, as follows:
  \begin{itemize}[nosep]
  \item total distance (km) of the non-public experiments;
  \item path (km), in the case of repetitive paths;
  \item total time (h) in terms of continuous operation;
  \item time interval (day/week/month/year, or d/w/m/y) between the first and the last run.
  \end{itemize}
\item $\left[\textbf{DE11}\right]$ \textbf{Datasets} -- if and which public datasets are used in the experiments;
\item $\left[\textbf{DE12}\right]$ \textbf{Evaluation metrics} -- which metrics are used for evaluation.
\end{itemize}

In Section~\ref{sec:discussion:experiments}, a comparison table of the public datasets identified by the DE11 will contain the sensory setup, ground-truth data availability from the datasets, and the distance and time characteristics, similar to the data extraction items for non-public data, among other aspects. As a result, the distinction between public and non-public data availability represented in DE8, DE9, and DE10 allows to understand the distance and time characteristics of non-public data independently from the public datasets.

Although the data extraction phase in a systematic literature review usually does not remove any records, 33 of the analyzed 175 works have extended versions of the proposed methodologies, more detailed ones, or equivalent methods applied in different conditions.
Thus, these records are not included in the review to improve the discussion section in terms of singularity and originality of proposed approaches for the long-term localization and mapping problem.
The extracted information helped identifying the corresponding extended and more complete versions of these works.
A document containing the association of the removed versions to the records included in the review is available in the public GitHub repository, including their bibliographic references. 
Consequently, 142 original works are included in this review for an overview of these records in Section~\ref{sec:overview}, and their synthesis and discussion in Section~\ref{sec:discussion}. The information relative to the 12 data items for each of the included records is available in Appendix~\ref{a2:data-extraction} and also in the repository. The included works represent 34.55\% of the 411 eligible records for this review. This result indicates that the methodology followed in this review led to a high percentage of quality results.


\section{Results Overview}
\label{sec:overview}

In this section, the main goal is to overview the results not in terms of their scientific contribution but in terms of their bibliographic data for presenting an overview of the included records in the review. First, statistic results of the data sources in which the 144 included records could be identified in the methodology allow the evaluation of the coverage between the sources.
Next, the tool VOSviewer~\parencite{results:vosviewer:1,results:vosviewer:2} is used to obtain the co-occurrence analysis for the keywords and the authors. The former focus on the keywords recency and their occurrence in the sources, while the latter discusses the research networks between the authors, and the ones with more publications in long-term localization and mapping.
Lastly, two analysis are presented relative to the evolution of the publication year and most relevant publication venues.

\subsection{Data source}
\label{sec:overview:db}

The results on the identification phase are exported to BibTeX files from each data source. This exportation considers all the information available in the data sources, such as citation (e.g., author, title, publication venue, and type of record) and bibliographic (e.g., affiliation and the publisher) information of each record, the abstract, and author and indexed keywords. Next, using the \texttt{bibtexparser}\footnote{\url{https://bibtexparser.readthedocs.io/en/master/}} Python library, the BibTeX files are processed to identify uncompleted records. For example, the DOI must be specified and, if not available, the record's information must be manually completed with a corresponding URL. Then, considering the 144 included records in this review, a Python script searches each record in the BibTeX files corresponding to each data source. This search uses the DOI, URL, and title data to identify if a data source had in its identification results the searched record. Given that these three fields can contain lower and upper letters, the respective strings must be compared only after converting them to lower cases. As a result, the number of identified records by each data source of the 144 included ones in the review are the following ones:

\begin{itemize}[nosep]
\item \citetitle{methodology:search:db:acm}: 25 records (17.4\%);
\item \citetitle{methodology:search:db:dimensions}: 85 records (59.0\%);
\item \citetitle{methodology:search:db:ieee-xplore}: 68 records (47.2\%);
\item \citetitle{methodology:search:db:inspec}: 104 records (72.2\%);
\item \citetitle{methodology:search:db:scopus}: 122 records (84.7\%);
\item \citetitle{methodology:search:db:wos}: 105 records (72.9\%).
\end{itemize}

The database \citetitle{methodology:search:db:scopus} is the source that identified the greatest number of included records. This result was expected given that \citetitle{methodology:search:db:scopus} is considered as one of the largest curated databases~\parencite{methodology:search:db:coverage:dim-scopus-wos}, indexing more than 25000 active titles (e.g., conferences proceedings, journals) and 7000 publishers\footnote{\url{https://www.elsevier.com/solutions/scopus/how-scopus-works}}.
Two other sources with more than 70\% of identified records are \citetitle{methodology:search:db:inspec} and \citetitle{methodology:search:db:wos}. Similarly to \citetitle{methodology:search:db:scopus}, these two databases index also records from thousands of journals, conferences, and publishers\footnote{\url{https://www.elsevier.com/solutions/engineering-village/content/inspec}}\textsuperscript{,}\footnote{\url{https://clarivate.com/webofsciencegroup/solutions/web-of-science/}}.
Although \citetitle{methodology:search:db:dimensions} is also a bibliographic database covering millions of publications from thousands of sources, this database is the newest one (created in 2018) relative to the other three considered in this review (\citetitle{methodology:search:db:inspec}, \citetitle{methodology:search:db:scopus}, and \citetitle{methodology:search:db:wos}) and could be a factor to why it obtained a lower percentage (59.0\%) than the other three databases. Another possible reason is that \citetitle{methodology:search:db:scopus} and \citetitle{methodology:search:db:wos} have the majority of their coverage in Life Sciences, Physical Sciences, and Technology Area (including the Engineering subject area related to the topic of this review), while \citetitle{methodology:search:db:dimensions} has better coverage in Social Sciences and Arts \& Humanities~\parencite{methodology:search:db:coverage:dim-scopus-wos}.
Even though \citetitle{methodology:search:db:ieee-xplore} is a digital library and only indexes works published by IEEE and its partners, this data source returns 47.2\% of the include records in the review. The main reason is that this library indexes publications related to electrical engineering and computer science, subject areas related to long-term localization and mapping\footnote{\url{https://ieeexplore.ieee.org/Xplorehelp/overview-of-ieee-xplore/about-ieee-xplore}}.
Finally, the \citetitle{methodology:search:db:acm} using \textit{The ACM Guide to Computing Literature} collection only finds published records by ACM and possible links to other records focused exclusively on computing\footnote{\url{https://libraries.acm.org/digital-library/acm-guide-to-computing-literature}} and not directly related to the Computer Science or Engineering subject areas, explaining why this source obtained a lower coverage percentage of the included results than the other sources for this review.

Furthermore, Table~\ref{tab:overview:source} presents a coverage analysis of the identified results from each data source for the 144 included records in this review. Table~\ref{tab:overview:source:intersect} presents the pairwise overlap between sources. The corresponding percentage is the ratio of records identified by both sources to the one between the two that has the smallest number of results: $\#\{A\cap B\} / \text{min}\{\#A,\#B\}$, where $\#A$ and $\#B$ is the number of results for a data source $A$ and $B$, respectively, and $\#\{A\cap B\}$ is the intersection results between the two sources. For example, if the pairwise results is 100\%, it means that the data source with more records found was capable of obtaining all the results, i.e., had full coverage over the other source. Table~\ref{tab:overview:source:union} reports the percentage of records identified by at least one of two data sources over all 144 included records: $\#\{A\cup B\} / 144$, where $A\cup B$ is the union correspondence results of the sources $A$ and $B$. This percentage represents the joint coverage of two databases over the 144 included records.
					
\begin{table}[h]
  \centering
  \caption{Pairwise coverage analysis of the data sources considered in the review over the 144 included records: (a) identification only on both pairwise sources ($\#\{A\cap B\} / \text{min}\{\#A,\#B\}$); (b) on either ones ($\#\{A\cup B\} / \#\text{records}$). Legend: dim -- \citetitle{methodology:search:db:dimensions}, ieee -- \citetitle{methodology:search:db:ieee-xplore}, insp -- \citetitle{methodology:search:db:inspec}, scop -- \citetitle{methodology:search:db:scopus}, wos -- \citetitle{methodology:search:db:wos}.}
  \label{tab:overview:source}
  \subfloat[][]{%
  \label{tab:overview:source:intersect}%
  \begin{tabular}{l|P{0.08\columnwidth}P{0.08\columnwidth}P{0.08\columnwidth}P{0.08\columnwidth}P{0.08\columnwidth}P{0.08\columnwidth}}
\hline
$A\cap B$ & \textbf{acm} & \textbf{dim} & \textbf{ieee} & \textbf{insp} & \textbf{scop} & \textbf{wos}\\
\hline
\textbf{acm}  & -- & 96.0\% & 44.0\% & 88.0\% & 96.0\% & 96.0\%\\
\textbf{dim}  & -- & --     & 69.1\% & 77.6\% & 97.6\% & 95.3\%\\
\textbf{ieee} & -- & --     & --     & 89.7\% & 91.2\% & 73.5\%\\
\textbf{insp} & -- & --     & --     & --     & 87.5\% & 68.3\%\\
\textbf{scop} & -- & --     & --     & --     & --     & 89.5\%\\
\textbf{wos}  & -- & --     & --     & --     & --     & --\\
\hline
    \end{tabular}%
  }
  \linebreak
  \subfloat[][]{%
  \label{tab:overview:source:union}%
  \vspace{0.5em}
  \begin{tabular}{l|P{0.08\columnwidth}P{0.08\columnwidth}P{0.08\columnwidth}P{0.08\columnwidth}P{0.08\columnwidth}P{0.08\columnwidth}}
\hline
$A\cup B$ & \textbf{acm} & \textbf{dim} & \textbf{ieee} & \textbf{insp} & \textbf{scop} & \textbf{wos}\\
\hline
\textbf{acm}  & -- & 59.7\% & 56.9\% & 74.3\% & 85.4\% & 73.6\%\\
\textbf{dim}  & -- & --     & 73.6\% & 85.4\% & 86.1\% & 75.7\%\\
\textbf{ieee} & -- & --     & --     & 77.1\% & 88.9\% & 85.4\%\\
\textbf{insp} & -- & --     & --     & --     & 93.8\% & 95.8\%\\
\textbf{scop} & -- & --     & --     & --     & --     & 92.4\%\\
\textbf{wos}  & -- & --     & --     & --     & --     & -\\
\hline
  \end{tabular}%
  }
\end{table}

Analyzing the coverage results in Table~\ref{tab:overview:source}, the first observation is that the pairwise union results of two sources increase the independent coverage of each source. This observation validates the need identified in the methodology discussed in Section~\ref{sec:methodology} to consider several data sources in the identification phase of a review. Moreover, the pairwise union coverage of \citetitle{methodology:search:db:inspec}, \citetitle{methodology:search:db:scopus}, and \citetitle{methodology:search:db:wos} is greater than 90\% of the included records. When evaluating the joint coverage of these three databases, they identify all 144 of the included records, i.e., a 100\% coverage. Although this result could indicate that those three sources guarantee full coverage of the long-term localization and mapping research topic, it is always advisable to consider as most as possible sources in the methodology. Another observation is relative to the overlap of \citetitle{methodology:search:db:scopus} with the other sources, which is greater than 85\%. This overlap indicates that \citetitle{methodology:search:db:scopus} covers results not only on the topic of this review but also the results obtained by the other sources considered in the methodology. Lastly, \citetitle{methodology:search:db:inspec} and \citetitle{methodology:search:db:wos} achieve a pairwise overlap percentage of 68.3\% between themselves, while their union represents 95.8\% of the included records. This discrepancy indicates that these two sources identify unique results between themselves. Indeed, \citetitle{methodology:search:db:inspec} identifies 33/144 records not found by \citetitle{methodology:search:db:wos}, and vice-versa for \citetitle{methodology:search:db:wos}, with 34/144 unique records.

\subsection{Keywords co-occurrence}
\label{sec:overview:kw}

Next, VOSviewer~\parencite{results:vosviewer:1,results:vosviewer:2} is used to analyze the co-occurrence of keywords in the included articles. This co-occurrence is the relatedness of items determined based on the number of documents in which the keywords occur together.
For this analysis, first, a Python script processes the BibTeX file containing the citation and bibliographic information, the author and the indexed keywords, and the abstract of the records to join the author with the indexed keywords in the same \texttt{keywords} field. Then, an online tool\footnote{\url{https://www.bibtex.com/c/bibtex-to-ris-converter/}} converts this processed BibTeX to a RIS file.
Even though VOSviewer supports file types directly exported from \citetitle{methodology:search:db:dimensions}, \citetitle{methodology:search:db:scopus}, or \citetitle{methodology:search:db:wos} as input, none of these data sources obtained all the 144 included records of the review in the identification phase. Given that VOSviewer does not support BibTeX files, the conversion to RIS file is required for using as input.
The disadvantage of using this file format in VOSviewer is only allowing to perform co-occurrence of items (e.g., keywords or authors), while bibliographic data from \citetitle{methodology:search:db:dimensions}, \citetitle{methodology:search:db:scopus}, or \citetitle{methodology:search:db:wos} in CSV files would allow other analysis such as citation, co-citation, or bibliographic coupling.
However, the creation of these CSV files follow different templates depending on the data source. So, RIS files allow the integration of all 144 included records for obtaining the two co-occurrence analysis presented in this review (namely, keywords and co-authorship).

In Figure~\ref{fig:overview:kw:original}, the network presents the overlay visualization of the keywords co-occurrence in the included records weighted by the number of occurrences of each term, using full counting for the links' strength. The latter computes the strength of the links directly by the number of co-occurrences of the respective two terms.
The overlay visualization colors the keywords differently according to the average publication year of the included records in which each of the keywords appears. This coloring allows analyzing which are the ones that are associated with the most recent publications.
As for the keywords' weighting, the number of occurrences dictates the size of the circles.
Furthermore, the minimum number of occurrences of a keywords set in VOSviewer for obtaining the graph is 5 originating the 35 keywords illustrated in Figure~\ref{fig:overview:kw:original}. This parameter was selected for visualization purposes while also filtering uninteresting keywords.
Similarly, setting the attraction and repulsion parameters to 2 and 0, respectively, distances the terms more from each other than using the values recommended in the VOSviewer manual\footnote{\url{https://www.vosviewer.com/documentation/Manual_VOSviewer_1.6.8.pdf}} (2 and 1, respectively). These two parameters only interfere in the localization of the terms in the map, not in the graph connections.
Lastly, a thesaurus of the keywords (available in the repository) is used to join similar terms: spelling differences (e.g., localization -- localisation), full terms versus abbreviations (simultaneous localization and mapping -- SLAM), while also allowing the concatenation of long keywords for visualization reasons.

\begin{figure*}[!t]
  \centering
  \subfloat[][]{\includegraphics[width=\columnwidth]{figures/kw.png}%
  \label{fig:overview:kw:original}}
  \hfill
  \subfloat[][]{\includegraphics[width=\columnwidth]{figures/kw_long-term.png}%
  \label{fig:overview:kw:long-term}}
  \caption{Keywords co-occurrence analysis on the 144 included records generated by VOSviewer with overlay visualization by the average publication year: (a) original keywords; (b) all keywords containing long-term and lifelong summarized by the terms themselves. Parameters used for generating the co-occurrence network: minimum number of occurrences = 5, attraction = 2, repulsion = 0, scale = 1.49, circles size variation = 0.5, lines size validation = 1.0. Legend: \texttt{cnn} -- Convolutional Neural Networks, \texttt{env} -- environment, \texttt{localiz} -- localization, \texttt{nav} -- navigation, \texttt{nn} -- Neural Networks, \texttt{recog} -- recognition, \texttt{vslam} -- visual SLAM.}
  \label{fig:overview:kw}
\end{figure*}

Overall, the keyword \texttt{robot} is the one that appears more times in the included records: 111 occurrences, links with 34 other terms, and has a total link strength of 403 (sum of co-occurrences of all of its links). This result is expected due to the relation of this review's topic to robotics.
Similarly, three other keywords in the network related to long-term localization and mapping topic with high values of occurrence, number of links, and total link strength are \texttt{slam} (75, 34, and 288), \texttt{mapping} (48, 33, and 204), and \texttt{localization} (47, 32, and 194, respectively). The methodology for the search strategy discussed in Section~\ref{sec:methodology:search} considers all of these four keywords. Thus, the significant influence of \texttt{robot}, \texttt{slam}, \texttt{mapping}, and \texttt{localization} in the keywords co-occurrence analysis indicates that, after the all the phases executed in this review's methodology, the 144 included records have a high correlation with the keywords considered in the search query. Given that the keywords are usually selected or indexed to capture the essence of the document, this correlation indicates that the search query is appropriate to obtain the search results, even considering only the keywords as search fields.

As for keywords related to the outcome of the PICO framework, \texttt{long-term autonomy} occurs only 6 times in the included records, linking with 16 other keywords and having a total link strength of 27. This low occurrence could indicate that the term \texttt{long-term autonomy} is not usually used by the authors nor indexed by the databases. 
However, the specific term of \texttt{long-term autonomy} does not summarize all the possibilities for the outcome of the PICO framework (see Section~\ref{sec:purpose}). Indeed, for this reason, the search query for the identification phase uses only the following single terms: \texttt{"long term"} and \texttt{"life long"} (resumes the possibility of having a space or a hyphen), and \texttt{lifelong}.
Figure~\ref{fig:overview:kw:long-term} presents the keywords co-occurrence analysis using the same parameters for obtaining Figure~\ref{fig:overview:kw:original}. The difference to the latter network is using a thesaurus that summarizes all the keywords that contain \texttt{long-term} and \texttt{lifelong} into the terms themselves, obtaining 36 keywords with a minimum of 5 occurrences in the 144 included records.
In terms of occurrences, number of links, and total link strength, the impact of the thesaurus keyword \texttt{long-term} is 25, 28, and 105, and for lifelong 6, 17, and 31, respectively. These values are much higher than the ones respective only to \texttt{long-term autonomy} from Figure~\ref{fig:overview:kw:original}.
The reason is that \texttt{long-term} in Figure~\ref{fig:overview:kw:long-term} compiles the occurrences of keywords such as \texttt{long-term autonomy}, \texttt{long-term mapping}, and \texttt{long-term localization} (6, 2, and 2 occurrences, respectively), and \texttt{lifelong} sum up, for example, three different versions of \texttt{lifelong learning} (using \texttt{lifelong}, \texttt{life-long} and \texttt{life long} with 2, 1, and 2 occurrences, respectively) and \texttt{lifelong slam} (1 occurrence). Hence, these results proves that the third \texttt{AND} part of the search query (\texttt{"long term" OR "life long" OR lifelong}) covers well the PICO framework's outcome. Plus, they also show no consensus among the authors and by the databases indexation on how to define a keyword for the topic of long-term localization and mapping.

In terms of the average year of publication, analyzing the diagrams in Figure~\ref{fig:overview:kw} on its colorization, the first observation is the recency of terms related to visual localization. The keywords visual SLAM (\texttt{vslam}), visual navigation (\texttt{visual nav}), and visual localization (\texttt{visual localiz}) have all an average publication year higher than 2017. This recency indicates that recent approaches related to the topic of this review, long-term localization and mapping, are more inclined to use vision as a sensorization input.
Another sensor that appeared with high relevance in the network is \texttt{radar}, with 15 occurrences and an average publication year of 2019.20. This sensor is agnostic to the environment changes such as illumination and season changes intrinsically associated with vision and could be the reason why the recent works related to long-term localization and mapping are using it.
Moreover, place recognition (\texttt{place recog}) stands out not only by its recency but importance. The keyword itself (\texttt{place recog}) occurs 31 times and an average publication year of 2017.77, with terms related to place recognition such as \texttt{loop closure} and global localization (\texttt{global localiz}) with recent average publication years (2017.82 and 2018.75, respectively) and strong link to place recognition (5 co-occurrences for each of the links between \texttt{loop closure} and \texttt{global localiz} with \texttt{place recog}). Lastly, machine learning also seems to be used in recent works included in this review. The keywork learning occurs 7 times with an average publication year of 2017.00. Neural Netowrks (\texttt{nn}), Convolutional Neural Networks (\texttt{cnn}), and \texttt{deep learning} have a similar number of occurrences (6, 5, and 8) and publication years higher than 2017 (2017.83, 2018.00, and 2019.12, respectively). These results could mean another trend of using machine learning to improve the long-term autonomy of mobile robots.

Although the recency of keywords related to dynamic environments is lower than 2017 (2015.50 and 2016.75 for \texttt{dynamic env} and \texttt{env change}), they have a high occurrence (14 and 12, respectively), located close to each other in the network, and have a strong link between them (4 co-occurrences). Three keywords also located near each other are \texttt{graph}, \texttt{marginalization}, and information theory (\texttt{info theory}) while having similar average publication years (2014.36, 2014.00, and 2015.50, respectively). Even though the number of occurrences of these terms is low (11, 6, and 5 for \texttt{graph},\texttt{ marginalization}, and \texttt{info theory}, respectively), their map proximity could indicate a focus in the past on the topic of graph sparsity, i.e., maintaining the graph in the long-term to only depend on the environment size and not on the robot's operation time.

The keywords co-occurrence analysis also relates to the categories of DE1 (see Section~\ref{sec:methodology:data}). Works associated with place recognition, global localization, and loop closure terms require invariance to the appearance changes in the environment, equivalent to the appearance category. The dynamics category is associated with works focused on dynamic environments. As for the other group of keywords with a high occurrence and strong links between each other, the ones related to graph and information theory, the respective works focus on removing uninformative data from the map~\parencite{kretzschmar-stachniss:2012:0278364912455072}, which is related to map sparsification, and so, to the sparsity category of DE1. These relations between the appearance, dynamics, and sparsity categories to the semantic analysis of the keywords co-occurrence supports the categorization of DE1 considered in this review, while also indicating that the discussion on the proposed methodologies should focus on each one of the categories. Even though the two remaining categories of DE1 (multi-session and computational) are not represented in the keyword analysis, the execution of the data extraction phase identified the need for having these two categories, given the importance of multi-session handling and computational efficiency for long-term localization and mapping. However, each category of DE1 will be discussed in Section~\ref{sec:discussion} in further detail.


\subsection{Co-authorship analysis}
\label{sec:overview:authors}

The other analysis obtained using VOSviewer is the co-authorship network presented in Figure~\ref{fig:overview:authors}. Similar to the keywords network illustrated in Figure~\ref{fig:overview:kw}, the co-occurrence of the authors' names creates links among them in the graph. The strength of these links is dictated by the number of documents the two authors of a link are co-authors in the same record, and the number of co-authored works determines the size of the circles respective to each author in the graph. 
In contrast to Figure~\ref{fig:overview:kw}, the network in Figure~\ref{fig:overview:authors} does not have any overlay specific to coloring depending on the average publication year. Instead, the main goal of the co-authorship analysis in this review is to present possible research networks detected in the 144 included records. Thus, the coloring in Figure~\ref{fig:overview:authors} represents the clusters of authors detected by VOSviewer. This network only considers authors with a minimum of 3 works for relevance and visualization reasons, resulting in 29 authors.
Also, authors identified only by the initial of the first name and by the surname can lead to incorrect correspondences in terms of co-authorship. VOSviewer detects 392 authors in the 144 included records using the original RIS file used in Section~\ref{sec:overview:kw} compared to 413 after checking the authors names. Indeed, a manual check is performed on all authors of the included records to guarantee no false correspondences for the co-authorship analysis with VOSviewer. This manual check ensures each author has its full first and surname and any middle initials while also using the same name for an author in different records.

\begin{figure}[h]
  \centering
  \includegraphics[width=\columnwidth]{figures/authors.png}
  \caption{Co-authorship analysis on the 144 included records generated by VOSviewer. Parameters used for generating the co-occurrence network: minimum number of occurrences = 3, attraction = 2, repulsion = -3, scale = 1.49, circles size variation = 1.0, lines size validation = 1.0.}
  \label{fig:overview:authors}
\end{figure}

Analyzing Figure~\ref{fig:overview:authors}, the co-authorship network presents 8 clusters. These clusters are separated from each other, i.e., no link exists between authors from different clusters. However, this separation does not mean that there is not any co-authorship between authors from different clusters only indicating that for a miminum of 3 co-authored documents there is not a connection between these 8 clusters. Even so, the graph presented in Figure~\ref{fig:overview:authors} allows the identification of the most relevant research networks in terms of number of co-authored documents and in the context of long-term localization and mapping, considering the 144 records included in this review. As a results, the following enumeration presents the authors that belong to each cluster in the format of author (number of co-authored documents):

\begin{enumerate}\setlength\itemsep{-0.5em}
\item Rong Xiong \orcid{0000-0001-9318-9014}\textsuperscript{,}\scholar{1hI9bqUAAAAJ} (7),
        Yue Wang \orcid{0000-0002-0981-935X}\textsuperscript{,}\scholar{N543LSoAAAAJ} (7),
        Huan Yin \orcid{0000-0002-0872-8202}\textsuperscript{,}\scholar{1fNc3vUAAAAJ} (6),
        Li Tang \orcid{0000-0003-2590-6872} (5),
        and Xiaqing Ding \orcid{0000-0001-7802-0130}\textsuperscript{,}\scholar{6u5OHUcAAAAJ} (4);
\item Hao Zhang \scholar{Ug2VxyUAAAAJ} (6),
        John J. Leonard \orcid{0000-0002-8863-6550}\textsuperscript{,}\scholar{WPe7vWwAAAAJ} (5),
        Michael Kaess \scholar{27eupmsAAAAJ} (5),
        Fei Han \orcid{0000-0002-8619-3987} (3),
        Guoquan Huang \scholar{trMUyZIAAAAJ} (3),
        José Neira \orcid{0000-0003-0668-977X}\textsuperscript{,}\scholar{scoMbR8AAAAJ} (3),
        and Ryan M. Eustice \orcid{0000-0002-9989-4942}\textsuperscript{,}\scholar{WroYmiAAAAAJ} (3);
\item Wolfram Burgard \orcid{0000-0002-5680-6500}\textsuperscript{,}\scholar{zj6FavAAAAAJ} (7),
        Cyrill Stachniss \orcid{0000-0003-1173-6972}\textsuperscript{,}\scholar{8vib2lAAAAAJ} (5),
        Giorgio Grisetti \orcid{0000-0002-8038-9989}\textsuperscript{,}\scholar{yD-SFG4AAAAJ} (3),
        Henrik Kretzschmar \scholar{pZuaUlwAAAAJ} (3),
        and Tayyab Naseer \orcid{0000-0002-3350-3005}\textsuperscript{,}\scholar{1FePZqEAAAAJ} (3);
\item Roland Siegwart \orcid{0000-0002-2760-7983}\textsuperscript{,}\scholar{MDIyLnwAAAAJ} (6),
        Igor Gilitschenski \orcid{0000-0001-6426-365X}\textsuperscript{,}\scholar{Nuw1Y4oAAAAJ} (3),
        Marcin Dymczyk \orcid{0000-0003-3667-8764}\textsuperscript{,}\scholar{XYHy7U8AAAAJ} (3),
        Michael Bosse \scholar{eopb1VgAAAAJ} (3),
        and Paul Furgale \orcid{0000-0002-7367-1046}\textsuperscript{,}\scholar{RNDtSG8AAAAJ} (3);
\item Tom Duckett \orcid{0000-0003-2971-7905}\textsuperscript{,}\scholar{et1GU2EAAAAJ} (6),
        Cheng Zhao \orcid{0000-0001-8502-3233}\textsuperscript{,}\scholar{EAC-8m0AAAAJ} (3),
        and Tomáš Krajník \orcid{0000-0002-4408-7916}\textsuperscript{,}\scholar{Qv3nqgsAAAAJ} (3);
\item Gordon F. Wyeth \orcid{0000-0002-4996-3612}\textsuperscript{,}\scholar{yfXZfXEAAAAJ} (5)
        and Michael J. Milford \orcid{0000-0002-5162-1793}\textsuperscript{,}\scholar{TDSmCKgAAAAJ} (4);
\item Paul Newman \scholar{BtO5fTUAAAAJ} (4);
\item Timothy D. Barfoot \orcid{0000-0003-3899-631X}\textsuperscript{,}\scholar{N_vPIhoAAAAJ} (3).
\end{enumerate}

When analyzing the affiliations of the authors mentioned previously at the time of publication, all authors of the first cluster belonged to the State Key Laboratory of Industrial Control and Technology (SKLICT) and the Institute of Cyber-Systems and Control at Zhejiang University in China. Even though Huan Yin, Yue Wang, Xiaqing Ding, Li Tang, and Rong Xiong mention their affiliation to the Joint Centre for Robotics Research between Zhejiang University, China, and the University of Technology Sydney, Sydney, in the work \parencite{yin-et-al:2020:2905046}, this specific affiliation only appeared in this article. The total link strength (sum of all links weights) of each of the authors in that cluster is higher than 16, meaning a high co-authorship between them. Indeed, all five authors have links between all of them.
Similar to the first cluster, the third, fourth, fifth, and sixth clusters have common affiliations within each one: the Autonomous Intelligent Systems at the University of Freiburg in Germany, the Autonomous Systems Lab (ASL) at ETH Zürich in Switzerland, the Lincoln Centre for Autonomous Systems (LCAS) at the University of Lincoln in UK, and the School of Electrical Engineering and Computer Science at Queensland University of Technology (QUT) in Australia.
However, the interlinking between the authors is not as strong as in the first cluster, as shown in Figure~\ref{fig:overview:authors} by the authors of these clusters not being connected between all the ones within each cluster. Even so, the common affiliation shows there is considerable interest by these research units in the long-term localization and mapping topic.

The affiliation analysis in the second cluster is more complex given that there was no affiliation common to all authors at the time of the records' publication. Instead, the following affiliations were found: Fei Han and Hao Zhang with the Department of Computer Science at Colorado School of Mines in the USA, Guoquan Huang with the Department of Mechanical Engineering at the University of Delaware in the USA, John J. Leonard and Michael Kaess with the Computer Science and Artificial Intelligence Laboratory (CSAIL) at the Massachusetts Institute of Technology (MIT) in the USA, Ryan M. Eustice with the Perceptual Robotics Laboratory (PeRL) at the University of Michigan in the USA, and José Neira with the Instituto Universitario de Investigación en Ingeniería de Aragón (I3A) at the Universidad de Zaragoza in Spain.
Although there are 5 different affiliations to which the 7 authors stated in the respective records, 4 of the research institutions noted for the second cluster are in the USA, indicating a possible reason for facilitating the linkage between these authors from different research units.

In terms of the clusters composed by single authors, the affiliations of Paul Newman and Timothy D. Barfoot are the Oxford Robotics Institute at the University of Oxford in UK and the Autonomous Space Robotics Laboratory (ASRL) at the University of Toronto Institute for Aerospace Studies (UTIAS) in Canada, respectively. Even though these two authors are not linked with any others in the network, the co-authorship analysis indicates that they have an interest in long-term localization and mapping. This interest is shown by their number of co-authored records: 4 and 3 by Paul Newman and Timothy D. Barfoot, respectively.

As for the number of co-authored publications, considering the 144 included records, the authors that appeared to have more research on the review's topic are Rong Xiong, Yue Wang, and Wolfram Burgard, given the 7 co-authored publications of each one. However, Rong Xiong and Yue Wang have co-authored the 7 documents attributed to each of them. This relation and similar ones can biase the analysis of which authors are having more impact in the review's topic.
The clustering shown in Figure~\ref{fig:overview:authors} allows a more unbiased analysis relative to the co-authorship links between authors. Thus, based on the clustering and which author from each cluster has the most co-authored publications, the most influential authors in long-term localization and mapping are the following ones: Rong Xiong (or Yue Wang), Hao Zhang, Wolfram Burgard, Roland Siegwart, Tom Duckett, Gordon F. Wyeth, Paul Newman, and Timothy D. Barfoot.

\subsection{Year of publication}
\label{sec:overview:year}

The relevance of the long-term localization and mapping topic can be evaluated by the evolution of the number of publications. Figure~\ref{fig:overview:year} presents this evolution from the earliest year of publication of the included records to the year at the time of writing this article. The latter has its respective data dashed to indicate that the last year is not completed at the time of writing. Analyzing Figure~\ref{fig:overview:year}, this review's topic seems to have gain relevance in 2009 with 6 works, compared to only one publication in 2007 and another in 2002 in the previous years to 2009. From that year onwards, the graph has an almost linear tendency reaching a maximum of 23 records in 2021, while already having 8 publications in 2022 until May 17, 2022. This tendency shows that long-term localization and mapping is gaining interest throughout the years and, consequently, supports the importance and relevance of this review for the scientific community.

\begin{figure}[!t]
  \centering
  \includegraphics[width=\columnwidth]{figures/year.png}
  \caption{Evolution of published records per year considering the 144 included records in this review}
  \label{fig:overview:year}
\end{figure}

\subsection{Publication venue}
\label{sec:overview:publication}

Finally, the last overview of the 144 included records in the review is relative to the publication venue. Table~\ref{tab:overview:publication} presents the venues with more than 1 publication, separating the journals and conferences in two different tables (Tables~\ref{tab:overview:publication:journal} and \ref{tab:overview:publication:conference}, respectively).
The columns $\mu$ present the average year of publication of the records associated to a certain venue, while $\text{max}$ columns display the publishing recency by the year of the most recent publication in the venue. For comparing to the average value ($\mu$), the third column ($\sigma$) of each table presents the standard deviation based on the publication year data.
The last column state the number of records published in the venue from the 144 records included in the review for discussion.

\begin{table}[!h]
  \centering
  \caption{Publication venues of the included records in this review with more than one record published in the venue: (a) journals; (b) conferences. Legend: $\mu$ -- average year of publication, $\sigma$ -- standard deviation of the publication year, $\text{max}$ -- maximum year of publication, \# -- number of records published at a certain venue}
  \label{tab:overview:publication}
  \subfloat[][]{%
  \begin{tabular}{p{0.55\columnwidth}cccc}
\hline
                 & \multicolumn{3}{c}{\textbf{Year}} & \\
\cline{2-4}
\textbf{Journal} & $\mu$ & $\sigma$ & $\text{max}$ & \textbf{\#}\\
\hline
Robotics and Autonomous Systems & 2016 & 3.9 & 2021 & 13\\
IEEE Robotics and Automation Letters & 2019 & 1.7 & 2022 & 12\\
International Journal of Robotics Research & 2014 & 3.2 & 2022 & 11\\
Journal of Field Robotics & 2017 & 3.5 & 2022 & 8\\
Autonomous Robots & 2017 & 2.2 & 2020 & 7\\
IEEE Transactions on Intelligent Transportation Systems & 2021 & 0.8 & 2022 & 4\\
Sensors & 2019 & 0.8 & 2020 & 4\\
IEEE Transactions on Robotics & 2017 & 3.1 & 2022 & 4\\
IEEE Sensors Journal & 2020 & 1.5 & 2021 & 2\\
International Journal of Advanced Robotic Systems & 2020 & 1.5 & 2021 & 2\\
\hline
  \end{tabular}\label{tab:overview:publication:journal}%
  }
  \linebreak
  \subfloat[][]{%
  \begin{tabular}{p{0.55\columnwidth}cccc}
\hline
                    & \multicolumn{3}{c}{\textbf{Year}} & \\
\cline{2-4}
\textbf{Conference} & $\mu$ & $\sigma$ & $\text{max}$ & \textbf{\#}\\
\hline
IEEE International Conference on Robotics and Automation (ICRA) & 2016 & 3.9 & 2021 & 22\\
IEEE/RSJ International Conference on Intelligent Robots and Systems (IROS) & 2017 & 3.8 & 2021 & 18\\
IEEE International Conference on Robotics and Biomimetics (ROBIO) & 2019 & 2.1 & 2021 & 3\\
IEEE International Intelligent Transportation Systems Conference (ITSC) & 2018 & 2.4 & 2021 & 3\\
European Conference on Mobile Robots (ECMR) & 2014 & 0.9 & 2015 & 3\\
IEEE Intelligent Vehicles Symposium (IV) & 2019 & 0.5 & 2019 & 2\\
International Conference on 3D Vision (3DV) & 2018 & 1.5 & 2019 & 2\\
International Conference on Advanced Robotics (ICAR) & 2011 & 2.0 & 2013 & 2\\
\hline
  \end{tabular}\label{tab:overview:publication:conference}%
  }
\end{table}

In terms of journals, the Robotics and Autonomous Systems, IEEE Robotics and Automation Letters, and the International Journal of Robotics stand out with more than 10 publications. Also, these journals have a high standard deviation (greater than 1.5), indicating that the publications spread out throughout the years.
In the case of the IEEE Robotics and Automation Letters, these results gain more relevance indicating a recent trend on publishing on this journal, considering that its creation was only on 2015\footnote{\url{https://www.ieee-ras.org/publications/ra-l}}.
With more than 5 publications, the Journal of Field Robotics and the Autonomous Robots have recent average of publication (2017) with a high standard deviation (greater than 2.0), similarly indicating that authors have been publishing in these two journals along the years.
In contrast, the IEEE Transactions on Intelligent Transportation Systems and Sensors journals have a standard deviation lower than 1 year, with an average publication year of at least 2019. The recency of publication on these two journals with a very low deviation suggests a recent interest of the authors to publish in these two journals works related to long-term localization and mapping.

As for conferences, the data in Table~\ref{tab:overview:publication:conference} shows a high discrepancy in the number of publications related to this review's topic in ICRA and IROS compared to the other venues. Indeed, all the other conferences have only a maximum of 3 records published in them, compared to 22 and 18 papers in ICRA and IROS, respectively. When considering that 62 of the 144 included records are published in conferences, ICRA and IROS with a total of 40 published works related to this review's topic represent 65\% of works published in conferences and 27.8\% of all included records. This result expresses the high relevance of ICRA and IROS in the topic of long-term localization and mapping.


\section{Discussion}
\label{sec:discussion}

\subsection{Appearance variance}
\label{sec:discussion:appearance}

\subsection{Dynamic environments}
\label{sec:discussion:dynamics}

\subsection{Map sparsification}
\label{sec:discussion:sparsify}

\subsection{Multi-session}
\label{sec:discussion:multisession}

\subsection{Computational}
\label{sec:discussion:computational}

\subsection{Evaluation metrics}
\label{sec:discussion:metrics}

\subsection{Long-term experimental data}
\label{sec:discussion:experiments}

\onecolumn

\renewcommand{\arraystretch}{1.125}
\setlength{\tabcolsep}{1.5pt}

\begin{tiny}

\begin{longtable}{@{\extracolsep{1pt}}
p{0.065\textwidth}                          % Dataset name
c|                                          % Year
cccccc|                                     % Long-term
p{0.095\textwidth}                          % Environment
p{0.095\textwidth}|                         % Domain
c|                                          % Sensory setup: odo
ccccccc|                                    %                camera
cc|                                         %                laser
cccc|                                       %                radar, sonar, imu, gps
cc|                                         % Sensor calibrations
p{0.095\textwidth}p{0.095\textwidth}|       % Ground-truth + Format
ccccc@{}}                                   % Distance, time, \#sequences
  \caption{Datasets used in the 144 included records for long-term localization and / or mapping experiments. Legend: odo -- odometry (wheeled, laser, visual, inertial, or a combination of odometry sources, dist. -- total distance length of the dataset, path -- total path distance if repeated several times, time -- total operation time, int. -- time interval between the start and end acquisition dates / time instants (d/w/m/y equivalent to day/week/month/year, 0 if only 1 run), and seq. -- number of sequences of the dataset.}
  %\vspace{0.5em}
  \label{tab:discussion:datasets}\\

%% FIRST TABLE HEADER
\hline
&&
\multicolumn{6}{c|}{\textbf{Long-term}} &
&&
\multicolumn{14}{c|}{\textbf{Sensor}} &
\multicolumn{2}{c|}{\textbf{Calib.}} &
&&
&&&&\\
&&&&&&&&&&
\multicolumn{1}{c}{} &
\multicolumn{7}{c}{\textbf{camera}} &
\multicolumn{2}{c}{\textbf{laser}}  &
&&&&
&&
&&
&&&&\\
\cline{12-18}
\cline{19-20}
\textbf{Dataset} & \textbf{Year} &
\rotatebox{90}{\textbf{lighting}} & \rotatebox{90}{\textbf{day/night}} & \rotatebox{90}{\textbf{weather}} & \rotatebox{90}{\textbf{seasonal}} & \rotatebox{90}{\textbf{dynamics}} & \rotatebox{90}{\textbf{sparsity}} & 
\textbf{Environ.} & \textbf{Domain} & 
\multicolumn{1}{c}{\rotatebox{90}{\textbf{odo}}} &
\rotatebox{90}{\textbf{gray}} & \rotatebox{90}{\textbf{color}} & \rotatebox{90}{\textbf{monocular}} & \rotatebox{90}{\textbf{stereo}} & \rotatebox{90}{\textbf{omni}} & \rotatebox{90}{\textbf{RGBD}} & \multicolumn{1}{c}{\rotatebox{90}{\textbf{thermal}}} &
\rotatebox{90}{\textbf{2D}} & \multicolumn{1}{c}{\rotatebox{90}{\textbf{3D}}} &
\rotatebox{90}{\textbf{radar}} & \rotatebox{90}{\textbf{sonar}} & \rotatebox{90}{\textbf{IMU}} & \rotatebox{90}{\textbf{GPS}} &
\rotatebox{90}{\textbf{intrinsic}} & \rotatebox{90}{\textbf{extrinsic}} & 
\textbf{GT data} & \textbf{Format} & 
\rotatebox{90}{\textbf{dist.} (km)} & \rotatebox{90}{\textbf{path} (km)} & \rotatebox{90}{\textbf{time} (h)} & \rotatebox{90}{\textbf{int.} (d/w/m/y)} & \rotatebox{90}{\textbf{\#seq.}}\\
\hline
\endfirsthead

%% TABLE HEADER IN THE FOLLOWING PAGES
\multicolumn{27}{l}{\itshape{\tablename\ \thetable{}: continued from previous page}}
\vspace{0.5em}\\
\hline
&&
\multicolumn{6}{c|}{\textbf{Long-term}} &
&&
\multicolumn{14}{c|}{\textbf{Sensor}} &
\multicolumn{2}{c|}{\textbf{Calib.}} &
&&
&&&&\\
&&&&&&&&&&
\multicolumn{1}{c}{} &
\multicolumn{7}{c}{\textbf{camera}} &
\multicolumn{2}{c}{\textbf{laser}}  &
&&&&
&&
&&
&&&&\\
\cline{12-18}
\cline{19-20}
\textbf{Dataset} & \textbf{Year} &
\rotatebox{90}{\textbf{lighting}} & \rotatebox{90}{\textbf{day/night}} & \rotatebox{90}{\textbf{weather}} & \rotatebox{90}{\textbf{seasonal}} & \rotatebox{90}{\textbf{dynamics}} & \rotatebox{90}{\textbf{sparsity}} & 
\textbf{Environ.} & \textbf{Domain} & 
\multicolumn{1}{c}{\rotatebox{90}{\textbf{odo}}} &
\rotatebox{90}{\textbf{gray}} & \rotatebox{90}{\textbf{color}} & \rotatebox{90}{\textbf{monocular}} & \rotatebox{90}{\textbf{stereo}} & \rotatebox{90}{\textbf{omni}} & \rotatebox{90}{\textbf{RGBD}} & \multicolumn{1}{c}{\rotatebox{90}{\textbf{thermal}}} &
\rotatebox{90}{\textbf{2D}} & \multicolumn{1}{c}{\rotatebox{90}{\textbf{3D}}} &
\rotatebox{90}{\textbf{radar}} & \rotatebox{90}{\textbf{sonar}} & \rotatebox{90}{\textbf{IMU}} & \rotatebox{90}{\textbf{GPS}} &
\rotatebox{90}{\textbf{intrinsic}} & \rotatebox{90}{\textbf{extrinsic}} & 
\textbf{GT data} & \textbf{Format} & 
\rotatebox{90}{\textbf{dist.} (km)} & \rotatebox{90}{\textbf{path} (km)} & \rotatebox{90}{\textbf{time} (h)} & \rotatebox{90}{\textbf{int.} (d/w/m/y)} & \rotatebox{90}{\textbf{\#seq.}}\\
\hline
\endhead

%% TABLE FOOTER
\hline
\endfoot
\hline
\endlastfoot



\citetitle{dataset:fhw} & 2001 &  &  &  &  &  & x & indoor (museum) & ground (TOURBOT) & x &  &  &  &  &  &  &  & x &  &  &  &  &  &  &  & -- & CARMEN & -- & -- & 1.98 & -- & 1\\
\hline
\citetitle{dataset:fr079} & 2003 &  &  &  &  &  & x & indoor (office) & ground (robot) & x &  &  &  &  &  &  &  & x &  &  &  &  &  &  &  & -- & CARMEN & -- & -- & 0.29 & -- & 1\\
\hline
\citetitle{dataset:fr101} & 2003 &  &  &  &  &  & x & indoor (office) & ground (robot) & x &  &  &  &  &  &  &  & x &  &  &  &  &  &  &  & -- & CARMEN & -- & -- & 0.29 & -- & 1\\
\hline
\citetitle{dataset:intel03} & 2003 &  &  &  &  &  & x & indoor (office) & ground (robot) & x &  &  &  &  &  &  &  & x &  &  &  &  &  &  &  & -- & CARMEN & 0.506 & -- & 0.75 & -- & 1\\
\hline
\citetitle{dataset:mit-kilian} & 2004 &  &  &  &  &  & x & indoor (office) & ground (robot) & x &  &  &  &  &  &  &  & x &  &  & x &  &  &  &  & -- & CARMEN & 2.2 & -- & 2.5 & -- & 1\\
\hline
\citetitle{dataset:city-center-fabmap} & 2008 & x &  &  &  & x &  & outdoor (urban) & ground (robot) &  &  & x & x &  &  &  &  &  &  &  &  &  & x & x &  & GPS, manual & plain text (non-img), jpg (img) & 2 & -- & -- & -- & 1\\
\hline
\citetitle{dataset:lip6indoor} & 2008 &  &  &  &  &  &  & indoor (office) & ground (handheld) &  &  & x & x &  &  &  &  &  &  &  &  &  &  & x &  & manual & ppm (imgs) & -- & -- & 0.11 & -- & 1\\
\hline
\citetitle{dataset:lip6outdoor} & 2008 & x &  &  &  & x &  & outdoor (campus) & ground (handheld) &  &  & x & x &  &  &  &  &  &  &  &  &  &  & x &  & manual & ppm (imgs) & -- & -- & 0.3 & -- & 1\\
\hline
\citetitle{dataset:new-college-fabmap} & 2008 & x &  &  &  & x &  & outdoor (campus) & ground (robot) &  &  & x & x &  &  &  &  &  &  &  &  &  & x & x &  & GPS, manual & plain text (non-img), jpg (img) & 1.9 & -- & -- & -- & 1\\
\hline
\citetitle{dataset:stlucia07} & 2008 & x &  &  &  & x &  & outdoor (urban) & ground (car) &  & x &  & x &  &  &  &  &  &  &  &  &  &  &  &  & -- & -- & 66 & -- & 1.67 & -- & 1\\
\hline
\citetitle{dataset:bicocca-indoor} & 2009 & x &  &  &  & x & x & indoor (office) & ground (Robocom) & x & x & x & x & x & x &  &  & x &  &  & x & x &  & x & x & map model, laser-based & plain text (non-img), png (img) & -- & -- & 2.5 & 3d & 5\\
\hline
\citetitle{dataset:cold} & 2009 & x & x & x &  & x &  & indoor (office) & ground (Pioneer 3, ATRV Mini, PeopleBot) & x &  & x & x &  & x &  &  & x &  &  &  &  &  &  &  & laser-based, manual & plain text (non-img), jpg (img) & 0.92 & -- & 0.99 & -- & 76\\
\hline
\citetitle{dataset:malaga-2009} & 2009 & x &  &  &  & x &  & outdoor (parking, campus) & ground (car) &  &  & x & x &  &  &  &  & x &  &  &  & x & x & x & x & RTK-GPS & Rawlog MRPT & 6.358 & -- & -- & -- & 6\\
\hline
\citetitle{dataset:newcollege} & 2009 & x &  &  &  & x &  & outdoor (campus) & ground (Segway) & x & x & x &  & x & x &  &  & x &  &  &  & x & x & x & x & GPS & plain text (non-img), png, jpg (img) & 2.2 & -- & 0.73 & -- & 1\\
\hline
\citetitle{dataset:albert-b} & 2010 & x &  &  &  &  & x & indoor (office) & ground (iRobot B21r) & x &  & x & x &  &  &  &  & x &  &  &  &  &  &  &  & -- & CARMEN (non-img), jpg (img) & -- & -- & 0.18 & -- & 1\\
\hline
\citetitle{dataset:cmu-vl} & 2011 & x &  &  & x & x &  & outdoor (urban) & ground (car) &  &  & x & x &  &  &  &  &  &  &  &  &  & x & x & x & GPS & -- & -- & 8.5 & -- & 1y & 16\\
\hline
\citetitle{dataset:ford} & 2011 & x &  &  &  & x &  & outdoor (campus, urban) & ground (car) &  &  & x &  &  & x &  &  & x & x &  &  & x & x & x & x & RTK-GPS & LCM log & -- & -- & -- & 2m & --\\
\hline
\citetitle{dataset:utias} & 2011 & x &  &  &  &  &  & indoor (empty space) & ground (iRobot Create) & x &  & x & x &  &  &  &  &  &  &  &  &  &  & x &  & external tracking system & jpg (img), dat (non-img) & -- & -- & 4.78 & -- & 9\\
\hline
\citetitle{dataset:alderley} & 2012 & x &  & x &  & x &  & outdoor (urban) & ground (car) &  &  & x & x &  &  &  &  &  &  &  &  &  &  &  &  & manual & -- & 16 & 8 & -- & -- & 2\\
\hline
\citetitle{dataset:tum-rgbd} & 2012 & x &  &  &  & x &  & indoor (office, industrial hall) & ground (handheld, Pioneer 3) &  &  & x &  &  &  & x &  &  &  &  &  & x &  & x & x & external tracking system & plain text (non-img), png (img + depth), ROS bag & 0.285 & -- & 0.35 & -- & 15\\
\hline
\citetitle{dataset:cobot} & 2013 & x &  &  & x & x & x & indoor (office) & ground (robot) & x &  & x &  &  &  & x &  & x &  &  &  &  &  &  & x & -- & ROS bag & 131 & -- & 260 & 2y3m & 1082\\
\hline
\citetitle{dataset:kitti} & 2013 & x &  &  &  & x &  & outdoor (urban) & ground (car) &  & x & x & x & x &  &  &  &  & x &  &  & x & x & x & x & RTK-GPS & png (img), binary (laser), plain text (imu, gps) & -- & -- & 1.18 & 8d & 61\\
\hline
\citetitle{dataset:mit-stata} & 2013 & x &  &  & x & x & x & indoor (office) & ground (PR2) & x &  & x &  & x &  & x &  & x &  &  &  & x &  &  & x & map model & ROS bag & 42 & -- & 38 & 1y9m & 84\\
\hline
\citetitle{dataset:nordlandsbanen} & 2013 & x &  & x & x &  &  & outdoor (railway) & ground (train) &  &  & x & x &  &  &  &  &  &  &  &  &  & x &  &  & GPS & mp4 (video stream), plain text (gps) & 2916 & 729 & 39.74 & -- & 4\\
\hline
\citetitle{dataset:gardens-qut} & 2014 & x & x &  &  & x &  & indoor, outdoor (campus) & ground (handheld) &  &  & x & x &  &  &  &  &  &  &  &  &  &  &  &  & ground-plane position & png (imgs), plain text (ground plane) & -- & -- & -- & -- & 3\\
\hline
\citetitle{dataset:lcas-strands} & 2014 & x & x &  & x & x &  & indoor (office) & ground (SCITOS-G5) &  &  & x &  &  &  & x &  & x &  &  &  &  &  &  & x & -- & ROS bag & -- & -- & -- & 1y1m & 368\\
\hline
\citetitle{dataset:kaist} & 2015 & x & x &  &  & x &  & outdoor (urban) & ground (car) &  &  & x &  & x &  &  & x &  &  &  &  & x & x & x & x & RTK-GPS & png (imgs), plain text (imu, gps) & 84 & -- & -- & 18d & 36\\
\hline
\citetitle{dataset:euroc} & 2016 & x &  &  &  &  &  & indoor (industrial hall, office) & air (AscTec Firefly) &  & x &  & x & x &  &  &  &  &  &  &  & x &  & x & x & external tracking system & ROS bag & 0.8936 & -- & 0.37 & -- & 11\\
\hline
\citetitle{dataset:nclt} & 2016 & x &  & x & x & x &  & indoor, outdoor (campus) & ground (Segway) &  &  & x &  &  & x &  &  & x & x &  &  & x & x & x & x & RTK-GPS, SLAM-based & binary (laser), tiff (img), plain text (non-laser or img) & 147.4 & -- & 34.9 & 1y4m & 27\\
\hline
\citetitle{dataset:kudamm} & 2017 & x &  &  &  & x &  & outdoor (urban) & ground (car) &  &  & x & x &  &  &  &  &  &  &  &  &  &  &  &  & manual & jpg (img) & -- & -- & -- & -- & 2\\
\hline
\citetitle{dataset:oxford-robotcar} & 2017 & x & x & x & x & x &  & outdoor (urban) & ground (car) & x &  & x & x & x &  &  &  & x & x &  &  & x & x & x & x & RTK-GPS & png (img), binary (laser), plain text (imu, gps, odo) & 1010.46 & 10 & -- & 1y8m & 133\\
\hline
\citetitle{dataset:yq21} & 2017 & x &  &  &  & x &  & outdoor (campus) & ground (car) &  &  & x & x & x &  &  &  &  & x &  &  & x & x & x & x & RTK-GPS & binary (laser), jpg (imgs), plain text (gps) & 23 & -- & 6.5 & 1w & 21\\
\hline
\citetitle{dataset:cmu-seasons} & 2018 & x &  &  & x & x &  & outdoor (urban) & ground (car) &  &  & x & x &  &  &  &  &  &  &  &  &  &  & x & x & manual & jpg (img) & -- & 8.5 & -- & 330d & 17\\
\hline
\citetitle{dataset:fas} & 2018 & x &  &  & x & x &  & outdoor (urban) & ground (car) &  &  & x &  & x &  &  &  &  &  &  &  &  & x &  &  & GPS, manual & jpg (img) & 110 & -- & -- & 3y & 3\\
\hline
\citetitle{dataset:robotcar-seasons} & 2018 & x &  & x &  & x &  & outdoor (urban) & ground (car) &  &  & x & x & x &  &  &  &  &  &  &  &  &  & x & x & manual & jpg (img) & -- & 10 & -- & 178d & 10\\
\hline
\citetitle{dataset:bonn} & 2019 &  &  &  &  & x &  & indoor (office) & ground &  &  & x &  &  &  & x &  &  &  &  &  &  &  & x & x & external tracking system & png (imgs, depth), plain text (imu, gps) & -- & -- & -- & -- & 26\\
\hline
\citetitle{dataset:cbd} & 2019 & x &  &  &  & x &  & outdoor (urban) & ground &  &  & x & x & x &  &  &  &  &  &  &  &  &  &  &  & manual & png (imgs) & -- & -- & -- & -- & 1\\
\hline
\citetitle{dataset:mulran} & 2020 &  &  &  &  & x &  & outdoor (urban) & ground (car) &  &  &  &  &  &  &  &  &  & x & x &  & x & x &  & x & SLAM-based & binary (laser), CSV (global pose, radar ray), png (radar polar img) & 41.2 & -- & -- & 2m15d & 12\\
\hline
\citetitle{dataset:robotcar-radar} & 2020 & x &  & x &  & x &  & outdoor (urban) & ground (car) & x &  & x & x & x &  &  &  &  & x & x &  & x & x & x & x & RTK-GPS, SLAM-based & png (img, raw laser, radar), binary (laser), plain text (imu, gps, odo) & 280 & 10 & -- & 1m & 32\\
\hline
\citetitle{dataset:usyd} & 2020 & x &  & x & x & x &  & outdoor (campus) & ground (car) & x &  & x & x &  &  &  &  &  & x &  &  & x & x & x & x & GPS & ROS bag & -- & -- & -- & 1y & 52\\
\hline
\citetitle{dataset:iplt} & 2021 & x & x & x & x & x &  & outdoor (parking) & ground (car) & x & x &  & x &  &  &  &  & x &  &  &  &  & x & x & x & GPS & ROS bag & -- & 0.2 & -- & 2y & 127\\
\hline
\citetitle{dataset:radiate} & 2021 & x & x & x &  & x &  & outdoor (parking, urban) & ground (car) &  &  & x & x & x &  &  &  &  & x & x &  & x & x & x & x & RTK-GPS & ROS bag & -- & -- & 4.98 & -- & --\\
\hline
\citetitle{dataset:ntu} & 2022 & x &  &  &  &  &  & indoor, outdoor (campus) & air (DJI M600) &  & x &  & x & x &  &  &  &  & x &  &  & x &  & x & x & external tracking system & ROS bag & 1.845 & -- & 0.9 & -- & 9\\



%  \end{tabular}}
%\end{table*}
\end{longtable}

\end{tiny}

\twocolumn

\renewcommand{\arraystretch}{1.25}
\setlength{\tabcolsep}{3pt}


\subsection{Final observations}
\label{sec:discussion:observations}


\section{Challenges and Future Directions}
\label{sec:future}

\begin{itemize}[nosep]
\item multi-robot long-term localization and mapping
\item fusion of different sensors
\item availability of datasets in continuous operation in challenging long-term scenarios (urban, industry -- storage / logistics facilities)
\item edge/cloud computing
\item parameters of the algorithms adapting along time accordingly to changes in the environment
\item specific hardware for improving computational performance (FPGA, operating systems specific for real-time usage, etc.)
\item interaction with the user: map update and consequent visualization by the users not being on raw data but instead on higher levels of abstraction (CAD, automatic extract of the environment topological locations)
\item hierarchical localization and mapping (topological location > coarse localization > millimeter level of localization or  something like that)
\end{itemize}

\section{Limitations of the Study}
\label{sec:limitations}

Section to discuss possible limitations of the study (timeframe, wider approach, etc.).

\begin{itemize}[nosep]
\item only one query for discussion: e.g., for searching datasets, possibly, a different query should have been used
\item overview long-term SLAM vs in-depth analysis and discussion of each type of techniques: our review synthesizes all types of techniques, if the reader wants an in-depth analysis, different reviews should be performed
\item related to previous one, each category appearance dynamics sparsity multi-session and computational have all different aspects related to each other, that probably should be treated differently in different data extraction items to improve the data organization -- however the main goal of this review was to overview all trends and not focus on one specifically, so this categorization should be done separately and probably in singular literature reviews.
\item limited information on the experiments conditions (traveled distance, duration, etc.) given by the authors
\item some works such as PointNetVLAD, FAB-MAP and SeqSLAM (+LOAM, and LEGO-LOAM) did not appeared in the identification phase. However, the ones included have improvements over all of these...; also ``Robust optimization of factor graphs by using condensed measurements''
\item discussion does not focus on the datasets used but rather on the methodologies, to not extend even further the review
\item ORBSLAM2, ORBSLAM3 not identified
\item in retrospective, a bias of the authors may occurred relative and influenced in the inclusion of articles related to air-vehicles, underwater, and multi-robot systems -- air although batteries possibly a limitation their viewpoint specially on bird eye camera view is completely different than usual ground robots (facing forward) and so long-term methods for ground may not be as effective due to different viewpoint, underwater due to low number of articles found and the importance of long-term SLAM in these environments (plants growing, different conditions of the water in camera), and multi-robot systems due to the future these systems can improve greatly the mapping and planning efficiency of robotic installations
\end{itemize}

\section{Conclusions}
\label{sec:conclusions}

This paper presents a systematic literature review on long-term localization and mapping for mobile robots. The review selects 142 works from the literature covering the main strategies to achieve lifelong SLAM and discussing the experimental data (including private experiments and public datasets) and evaluation metrics commonly used to assess the performance of autonomous systems in long-term operations.
The discussion analyzes the included works in terms of appearance invariance to changing conditions in the environment, dealing with dynamic elements in the scene, multi-session strategies, map sparsification techniques to bound the computational resources, and other computational-related topics.
Also, an overview over the bibliographic data of the included works identifies the most used terms in long-term SLAM and identifying research networks between authors using the VOSviewer~\parencite{results:vosviewer:1,results:vosviewer:2} tool, while also discussing the evolution of the number of publications over time and the publication venues with more records.

Overall, the methodologies discussed in this study are a step forward to achieve lifelong autonomy.
In terms of dealing with appearance changes in the environment, CNN-based features are more discriminative compared to handcrafted features, while considering both appearance and geometric cues in the descriptor improve its invariance to changing conditions. Semantic features and considering different sensorization sources (e.g., cameras, LiDAR, and/or radar) can also increase the robustness to appearance changes in the environment.
For dealing with dynamic elements, the most common approach is to distinguish between static, dynamic, and semi-static changes in the perceived environment, while only using static permanent changes for localization to improve the reliability of the pose estimator.
As for constraining the map size to the explored environment area instead of the trajectory length, the evaluation of the mutual information by using techniques based on information theory or heuristics is the most currently used technique to bound the map size.

However, there are still challenges and future directions for the research on long-term SLAM. Vision-based global recognition is affected by viewpoint variance, where using omnidirectional sensors to increase the field-of-view relative to perspective cameras may decrease that variance. Decentralized computation architectures may be an interesting research direction to offload heavier computation tasks from the mapping agent and improve the scalability of long-term SLAM algorithms. For example, cloud and edge computing devices allow the parallelization of tasks and increase the availability of computational resources, even though communication constraints should be accounted. Also, the extension of the methodologies discussed in this study to multi-robots systems may optimize the operation of autonomous systems while also increasing their robustness in the long-term.

Lastly, this review can be updated in future studies thanks to following a systematic review that defines explicitly the time interval in the selection process (the data of the last full inquiry is May 17, 2022) and the search query used to identify records from the literature.
Also, all documentation and scripts using during the review process are available in a public GitHub repository to facilitate replicating of the results and future updates to this review.
This paper may be extended to focus on single challenges of long-term SLAM for providing an in-depth comparison with experimental results between different methodologies.


\end{document}
